\documentclass[11pt,a4paper]{article}
\begin{document}

\newpage
\section*{개요: 중앙은행과 통화정책 모듈 3}

이 문서는 '중앙은행과 통화정책' 모듈 3의 핵심 내용을 담고 있으며, 통화정책이 경제에 미치는 영향을 심층적으로 분석합니다. 연방준비제도(Fed)의 이중 책무인 최대 고용과 물가 안정 달성을 위한 통화정책의 작동 방식, 인플레이션과 실업률 간의 상충 관계, 통화정책 전달 메커니즘, 그리고 통화정책의 한계점을 다룹니다. 특히 비전공자도 쉽게 이해할 수 있도록 각 개념에 대한 상세한 설명과 현실적인 예시를 제공하여 시험 대비 및 심화 학습에 도움을 줍니다.

\newpage
\section*{용어 정리}

다음 표는 본 문서에서 사용되는 주요 경제 및 통화정책 관련 용어들을 정리한 것입니다. 각 용어의 쉬운 설명과 원어를 함께 제시하여 이해를 돕습니다.

\begin{adjustbox}{width=\textwidth,center}
\begin{tabular}{lllc}
\toprule
\textbf{용어} & \textbf{쉬운 설명} & \textbf{원어} & \textbf{비고} \\
\midrule
\textbf{이중 책무} & 중앙은행의 두 가지 주요 목표: 물가 안정과 최대 고용 & Dual Mandate & Fed의 핵심 목표 \\
\textbf{GDP} & 한 국가에서 생산된 모든 재화와 서비스의 총 가치 & Gross Domestic Product & 경제 규모 측정 \\
\textbf{실질 GDP} & 인플레이션을 조정한 GDP & Real GDP & 실제 경제 성장률 반영 \\
\textbf{실업률} & 경제활동인구 중 실업자의 비율 & Unemployment Rate & 고용 상태 측정 \\
\textbf{자연 실업률} & 경제가 완전 고용 상태일 때도 존재하는 실업률 & Natural Rate of Unemployment & 마찰적/구조적 실업 포함 \\
\textbf{오쿤의 법칙} & 실업률과 GDP 성장률 간의 음의 상관관계 & Okun's Law & 통화정책의 지침 \\
\textbf{인플레이션} & 물가 수준의 지속적인 상승 & Inflation & 화폐 가치 하락 \\
\textbf{CPI} & 소비자가 구매하는 상품 및 서비스 가격의 변화 & Consumer Price Index & 소비자 물가 변동 측정 \\
\textbf{PPI} & 생산자가 판매하는 상품 및 서비스 가격의 변화 & Producer Price Index & 생산자 물가 변동 측정 \\
\textbf{필립스 곡선} & 인플레이션과 실업률 간의 상충 관계 & Phillips Curve & 단기적 정책 딜레마 \\
\textbf{인플레이션 기대} & 미래 인플레이션에 대한 경제 주체들의 예상 & Inflation Expectations & 필립스 곡선 이동 요인 \\
\textbf{생산량 갭} & 실제 GDP와 잠재 GDP의 차이 & Output Gap / GDP Gap & 경제 과열/침체 지표 \\
\textbf{인플레이션 갭} & 실제 인플레이션과 중앙은행 목표 인플레이션의 차이 & Inflation Gap & 물가 안정 목표 달성 여부 \\
\textbf{테일러 준칙} & 중앙은행이 금리를 설정하는 규칙 & Taylor Rule & 인플레이션 및 생산량 갭 반영 \\
\textbf{연방 기금 금리} & 은행 간 초단기 대출 금리 & Federal Funds Rate & Fed의 주요 정책 수단 \\
\textbf{순이자 마진} & 은행의 대출 이자 수입과 예금 이자 비용의 차이 & Net Interest Margin & 은행 수익성 지표 \\
\textbf{공급 충격} & 생산 비용에 갑작스러운 변화를 주는 사건 & Supply Shocks & 스태그플레이션 유발 가능 \\
\textbf{스태그플레이션} & 경기 침체와 높은 인플레이션이 동시에 발생하는 현상 & Stagflation & 통화정책의 큰 도전 \\
\bottomrule
\end{tabular}
\end{adjustbox}

\newpage
\section{모듈 3: 통화정책의 경제적 영향}

\subsection{레슨 3-0.1: 개요}

안녕하세요, 여러분. 어디에서 참여하시든 환영합니다. 연방준비제도(Federal Reserve, Fed)의 목적은 \textbf{최대 고용(maximum employment)}과 \textbf{물가 안정(stable prices)}을 달성하는 것입니다. 이를 \textbf{Fed의 이중 책무(dual mandate)}라고 부릅니다. 이상적인 세상에서는 경제가 최소한의 인플레이션을 겪고 모든 사람이 직업을 가질 수 있겠지만, 현실은 항상 그렇지 않습니다.

그렇다면 Fed의 통화정책 조치가 경제에 어떻게 영향을 미칠까요? 이것이 이 모듈의 목표입니다. 우리는 통화정책 조치와 가장 중요한 거시경제적 결과물들 사이의 연관성을 살펴볼 것입니다. 또한 이러한 변화가 어떻게 발생하는지도 논의할 것입니다.

\begin{conceptbox}
\textbf{Fed의 이중 책무 (Dual Mandate)}
\begin{itemize}
    \item \textbf{최대 고용 (Maximum Employment)}: 가능한 한 많은 사람들이 일자리를 가질 수 있도록 하는 것입니다. 단순히 실업률을 0으로 만드는 것이 아니라, 경제가 건강하게 작동할 때 자연적으로 발생하는 실업률 수준을 의미합니다.
    \item \textbf{물가 안정 (Stable Prices)}: 물가가 급격하게 오르거나 내리지 않고 안정적인 수준을 유지하는 것입니다. 이는 인플레이션이 너무 높거나 디플레이션이 발생하지 않도록 관리하는 것을 포함합니다.
\end{itemize}
이 두 가지 목표는 때때로 상충될 수 있어, Fed는 이 둘 사이에서 균형을 찾아야 합니다.
\end{conceptbox}

우리는 최대 고용이라는 아이디어부터 시작할 것입니다. 이중 책무는 일자리 측면에서 명시되어 있지만, 연방준비제도를 포함한 많은 중앙은행들은 \textbf{경제 성장(economic growth)}에 대해 많이 이야기합니다. 그 이유를 설명하기 위해 먼저 \textbf{국내총생산(Gross Domestic Product, GDP)}과 실업률 간의 관계를 살펴볼 것입니다.

그 다음, Fed의 목표 금리 인하가 어떻게 GDP를 증가시키고 결과적으로 실업률을 감소시킬 수 있는지 알아볼 것입니다. 이 논의의 일환으로, 미국 맥락에서 최대 고용의 목표가 무엇을 의미하는지 살펴볼 것입니다.

다음으로, 실업률, 임금, 인플레이션 간의 연관성을 연구할 것입니다. 실업률과 인플레이션 사이에는 \textbf{상충 관계(trade-off)}가 있다는 것을 알게 될 것입니다. 통화정책을 완화함으로써 중앙은행은 적어도 단기적으로는 실업률을 낮출 수 있지만, 이는 더 높은 인플레이션이라는 대가를 치를 수 있습니다. 실업률과 인플레이션 간의 이러한 상충 관계는 연방준비제도의 의사 결정에 핵심적인 요소입니다. 우리는 이 상충 관계가 유명한 \textbf{테일러 준칙(Taylor Rule)}과 같은 통화정책 규칙에 어떻게 내재되어 있는지 살펴볼 것입니다. 이러한 규칙은 경제 상황에 따라 금리가 정확히 어떻게 설정되어야 하는지를 지시합니다.

인플레이션과 실업률 간의 상충 관계를 염두에 두고, 통화정책이 어떻게 작동하는지, 즉 금리 변화가 경제 성장, 그리고 결과적으로 인플레이션과 실업률에 어떤 경로를 통해 영향을 미치는지 연구할 것입니다. 우리는 통화정책이 경제 전체에 미치는 영향을 보여주는 개괄적인 개요부터 시작할 것입니다. 그 다음, 통화정책이 은행의 신용 공급에, 그리고 결과적으로 경제의 통화 공급에 어떻게 영향을 미치는지 자세히 살펴볼 것입니다.

다음 레슨에서는 통화정책이 기업과 가계에 어떻게 영향을 미치고, 이것이 다시 투자 및 소비 결정에 어떤 영향을 미치는지 논의할 것입니다. 소비와 투자는 모두 GDP의 일부이며, 이러한 결정은 실업률에 영향을 미칠 것입니다. 통화정책은 경기 확장기에는 금리를 인상하고 경기 수축기에는 금리를 인하함으로써 경제 주기의 해로운 호황과 불황을 완화할 수 있습니다. 이것이 중앙은행이 달성하고자 하는 목표입니다.

마지막 두 레슨은 통화정책의 한계를 보여주는 사례 연구입니다. 우리는 통화정책이 유가 급등과 같은 공급 측면 충격(supply-side shocks)을 상쇄하는 데 거의 할 수 없다는 것을 알게 될 것입니다. 또한 구조적 경제 환경 변화로 인해 영구적으로 쇠퇴하는 산업의 고용을 지원하는 데도 많은 것을 할 수 없습니다.

요약하자면, 이전 모듈에서는 중앙은행이 사용할 수 있는 통화정책 도구와 그것이 금리에 어떻게 영향을 미치는지 논의했습니다. 이 모듈에서는 금리 변화가 경제에 어떻게 영향을 미치고, 이것이 연방준비제도가 물가 안정과 최대 고용이라는 이중 책무를 선택하는 데 어떻게 도움이 되는지 살펴봅니다.

\newpage
\section{레슨 3-1: 인플레이션과 실업률 간의 연관성}

\subsection{레슨 3-1.1: 경제 성장과 실업률 이해하기}

안녕하세요, 경제 성장과 실업률에 대한 강의에 오신 것을 환영합니다. 연방준비제도는 물가 안정과 최대 고용을 달성하기 위해 통화정책을 사용한다는 것을 기억하십시오. 하지만 최대 고용이란 정확히 무엇을 의미할까요? 이 수업에서는 실업률을 살펴볼 것입니다. 우리는 실업률에 대한 몇 가지 새로운 정의를 소개할 것입니다: \textbf{자연 실업률(natural rate of unemployment)}, \textbf{완전 고용(full employment)}, \textbf{최대 고용(maximum employment)}. 우리는 이 용어들이 무엇을 의미하고 왜 필요한지 이해할 것입니다. 또한 \textbf{국내총생산(Gross Domestic Product, GDP)}을 정의할 것입니다. 이것은 한 경제의 총 생산량을 측정하는 지표입니다. 우리는 데이터에서 실업률과 GDP 성장률 간의 관계를 살펴볼 것입니다. 마지막으로, 연방준비제도가 실업률을 0으로 줄이려고 하지 않는 이유를 이해할 것입니다.

\subsubsection*{실업률 데이터}

실업률 데이터부터 시작해 봅시다. 개인은 직업이 없고 구직 활동을 하고 있다면 실업 상태입니다. \textbf{실업률(unemployment rate)}은 경제활동인구 중 현재 실업 상태인 사람들의 비율입니다.

\begin{examplebox}
\textbf{실업률의 의미}
만약 한 국가의 경제활동인구가 100명이고, 이 중 5명이 직업이 없지만 적극적으로 일자리를 찾고 있다면, 실업률은 5\%가 됩니다. 나머지 95명은 고용 상태이거나, 아예 구직 활동을 하지 않는 비경제활동인구(예: 학생, 은퇴자)입니다.
\end{examplebox}

그래프를 보면 현대 시대 미국의 실업률을 알 수 있습니다. 그래프의 음영 처리된 부분은 \textbf{경기 침체(recessions)}를 나타냅니다. 경기 침체기에는 실업률이 급격히 증가하는 것을 볼 수 있습니다. 예를 들어, 1973년 10월부터 1975년 5월까지 실업률은 4.6\%에서 9\%로 거의 두 배가 되었습니다. 동시에, 경기 침체가 끝난 후에는 실업률이 빠르게 하락합니다. 1976년 3월에는 실업률이 이미 7.6\%로 떨어졌습니다.

\subsubsection*{경제 성장 측정: GDP}

이제 경제 성장의 측정 지표인 \textbf{국내총생산(Gross Domestic Product, GDP)}으로 넘어가겠습니다. GDP는 한 경제가 생산하는 모든 재화와 서비스의 총 달러 가치입니다. \textbf{실질 GDP(Real GDP)}는 이 측정치를 인플레이션에 대해 조정합니다. 미국의 실질 GDP는 시간이 지남에 따라 성장해 왔습니다. 그러나 경기 침체기에는 감소하며, 이는 경제가 수축한다는 것을 의미합니다. 데이터를 통해 이를 살펴보겠습니다.

그래프는 실질 GDP 성장률의 분기별 퍼센트 변화를 보여줍니다. 음영 처리된 경기 침체기에는 실질 GDP 성장률이 음수입니다. 1974년 각 분기 동안 GDP가 1\%만큼 감소한 것을 볼 수 있습니다.

\subsubsection*{실업률과 실질 GDP 성장률의 관계: 오쿤의 법칙}

실업률과 실질 GDP 성장률 간의 관계는 무엇일까요? 경기 침체기에는 실업률이 상승하고 GDP는 하락합니다. 경기 확장기에는 그 반대가 일어납니다. 이는 \textbf{음의 상관관계(negative correlation)}를 시사합니다.

이것이 어떻게 작동하는지 보기 위해 이전 두 차트를 결합해 봅시다. 각 점은 한 분기를 나타냅니다. X축에는 이전 분기 대비 실업률 증가율이 있고, Y축에는 해당 분기의 GDP 성장률이 있습니다. 데이터에서 음의 상관관계를 볼 수 있으며, 이는 아래로 기울어진 빨간색 선으로 표시됩니다. 역사적인 미국 데이터에 따르면, 이 선의 기울기는 약 -2입니다. 이 경험적 사실은 경제학자 아서 오쿤(Arthur Okun)에 의해 처음 언급되었으며, \textbf{오쿤의 법칙(Okun's Law)}으로 알려져 있습니다.

\begin{conceptbox}
\textbf{오쿤의 법칙 (Okun's Law)}
실업률과 실질 GDP 성장률 사이에는 음의 상관관계가 있다는 경제학적 원리입니다. 즉, 실업률이 1\% 감소하면 실질 GDP 성장률은 약 2\% 증가하는 경향이 있습니다 (미국 데이터 기준).

\textbf{왜 중요한가요?}
통화정책 입안자들은 오쿤의 법칙을 유용하게 활용합니다. 경제를 부양하기 위해 금리를 인하할 때 실업률이 얼마나 감소할지에 대한 지침을 제공하기 때문입니다. 금리 인하 후 실질 GDP가 얼마나 증가할지에 대한 추정치가 있다면, 실업률이 얼마나 감소할지도 알 수 있습니다. 이러한 거시경제 변수들 간의 관계를 이해하면 Fed가 통화정책을 더 잘 목표로 삼을 수 있습니다.
\end{conceptbox}

이제 다른 질문을 해봅시다.

\subsubsection*{자연 실업률: 실업률을 0으로 만들 수 없는 이유}

여러분은 중앙은행이 실업률을 0으로 줄일 수 없는 이유가 무엇인지 궁금할 수 있습니다. 답은 \textbf{자연 실업률(natural rate of unemployment)}이 존재하기 때문입니다. 이는 경제가 최고 성과를 내고 있을 때조차도 일부 근로자는 여전히 실업 상태에 있을 것이라는 의미입니다.

\begin{conceptbox}
\textbf{자연 실업률 (Natural Rate of Unemployment)}
경제가 완전 고용 상태에 있을 때도 존재하는 실업률입니다. 이는 경기 변동과 무관하게 발생하는 실업으로, 주로 다음과 같은 이유로 발생합니다.
\begin{itemize}
    \item \textbf{마찰적 실업 (Frictional Unemployment)}: 사람들이 새로운 직업을 찾거나, 교육을 마치고 첫 직장을 구하거나, 이사 등으로 인해 직업을 전환하는 데 시간이 걸리기 때문에 발생합니다.
    \item \textbf{구조적 실업 (Structural Unemployment)}: 기술 변화, 산업 구조 변화 등으로 인해 특정 기술을 가진 근로자들이 일자리를 잃고 새로운 기술을 배우거나 다른 산업으로 이동해야 할 때 발생합니다. (예: 자율 주행 기술이 트럭 운전사에게 미칠 영향)
    \item \textbf{제도적 요인}: 높은 실업 수당은 실업자들이 더 나은 일자리를 찾는 데 더 많은 시간을 할애하도록 유도할 수 있습니다.
\end{itemize}
따라서 중앙은행은 실업률을 0으로 만드는 것을 목표로 하지 않습니다. 대신 자연 실업률 수준에 도달하는 것을 '최대 고용'으로 간주합니다.
\end{conceptbox}

다시 미국의 실업률 시계열을 살펴보겠습니다. 대략적으로 말해, 우리는 좋은 경제 시기의 실업률을 기준으로 자연 실업률을 추정할 수 있습니다.

\begin{table}[h!]
\centering
\caption{미국 자연 실업률 추정치 (역사적 데이터)}
\label{tab:natural_unemployment_rates}
\begin{tabular}{lc}
\toprule
\textbf{기간} & \textbf{자연 실업률 추정치} \\
\midrule
1950년대 및 1960년대 & 4\% \\
1970년 ~ 1983년 & 6\% \\
1984년 ~ 2000년 & 5.5\% \\
\bottomrule
\end{tabular}
\end{table}

차트는 또한 실업률의 큰 주기적 변동을 보여줍니다. 특히 경기 침체기에는 실업률이 자연 실업률을 초과합니다. 바로 이 지점에서 통화정책이 효과적일 수 있습니다.

\subsubsection*{오늘날의 자연 실업률}

마지막 질문입니다. 오늘날의 자연 실업률은 얼마일까요? 새천년의 실업률을 살펴보겠습니다. 2001년 경기 침체 이후 실업률은 20세기 후반의 경기 침체에 비해 느리게 회복되었습니다. 마찬가지로, 2008년 금융 위기 이후 실업률은 예상만큼 빠르게 회복되지 않았습니다. 이는 자연 실업률이 증가했을 것이라는 추측을 불러일으켰습니다. 그러나 미국에서 가장 긴 경제 확장기가 끝나갈 무렵인 2019년 12월에는 실업률이 3.6\%였습니다. 마찬가지로, 코로나19 팬데믹 이후에도 실업률은 역사적으로 낮았습니다. 이는 자연 실업률이 감소했을 수 있음을 시사합니다. 전반적으로, 새로운 기술이나 노동 시장의 변화가 역할을 할 수 있지만, 우리가 결정하기까지는 더 많은 시간이 필요합니다.

\begin{summarybox}
\textbf{레슨 3-1.1 요약}
\begin{itemize}
    \item \textbf{오쿤의 법칙 (Okun's Law)}: 실업률 변화와 실질 GDP 변화는 데이터에서 음의 상관관계를 보입니다. 이 관계는 통화정책 결정에 중요한 지침이 됩니다.
    \item \textbf{자연 실업률 (Natural Rate of Unemployment)}: 경제가 완전한 생산 능력을 발휘할 때조차도 실업률은 0보다 클 것입니다. 이는 마찰적 실업과 구조적 실업 등 경제의 자연스러운 부분입니다.
\end{itemize}
\end{summarybox}

\newpage
\subsection{레슨 3-1.2: 실업률과 인플레이션 연결하기}

안녕하세요, 임금, 실업률, 인플레이션 간의 관계에 대한 강의에 오신 것을 환영합니다. \textbf{인플레이션 목표제(inflation targeting)}는 통화정책의 핵심입니다. 이 수업에서는 \textbf{노동 비용(labor costs)} 또는 \textbf{임금 인플레이션(wage inflation)}이 \textbf{소비자 물가 인플레이션(consumer price inflation)}으로 어떻게 이어지는지 논의할 것입니다. 우리는 임금과 인플레이션이 어떻게 함께 움직이는지 살펴볼 것입니다. 그 다음, 인플레이션과 실업률 간의 관계를 검토할 것입니다. 이를 위해 \textbf{필립스 곡선(Phillips Curve)}이라는 아이디어를 소개할 것입니다. 이것은 경제학에서 가장 유명한 차트 중 하나입니다. 인플레이션과 실업률 간의 상충 관계를 어떻게 설명하는지 알게 될 것입니다. 일반적으로 인플레이션이 낮아지면 실업률이 높아지는 대가를 치르게 됩니다. 마지막으로, \textbf{인플레이션 기대(inflation expectations)}의 변화가 필립스 곡선을 어떻게 이동시킬 수 있는지 이해할 것입니다.

\subsubsection*{소비자 물가 인플레이션의 원인}

경제학자, 정책 입안자, 언론은 항상 인플레이션에 대해 이야기합니다. 그들은 보통 \textbf{소비자 물가 인플레이션(consumer price inflation)} 또는 \textbf{CPI(Consumer Price Index)}를 언급합니다. 이는 소비자 상품 묶음의 가격 상승을 의미합니다.

여러분은 소비자 물가가 왜 오르는지 궁금할 수 있습니다. 한 가지 답은 \textbf{투입 가격(input prices)}입니다. 원유 가격이 높으면 자동차 연료비가 더 비싸다는 것을 우리 모두 알고 있습니다. 더 일반적으로, 투입 가격이 상승하면 상품 생산 및 서비스 제공 비용도 증가합니다. 다르게 말하면, \textbf{생산자 물가 인플레이션(producer price inflation)}이 높아지면 CPI가 증가할 수 있습니다. 이에 대한 데이터를 살펴보겠습니다.

그래프는 1980년대 이후 소비자 물가 인플레이션과 생산자 물가 인플레이션을 보여줍니다. 생산자 물가 지수(Producer Price Index, PPI)가 더 변동성이 크지만, 두 시리즈가 서로를 추적하는 것을 볼 수 있습니다. 생산자 물가가 높아지면 소비자 상품 가격도 높아집니다.

대부분 산업의 가장 큰 비용 구성 요소는 무엇일까요? 바로 \textbf{노동 비용(labor expenses)}입니다. 이 그래프에 \textbf{고용 비용 지수(employment cost index)}를 추가해 봅시다. 이 고용 비용 지수에는 임금과 복리후생이 모두 포함됩니다. 이 시리즈는 2000년에 시작되므로 시간 프레임을 제한해야 합니다. 우리는 다시 고용 비용과 소비자 물가 인플레이션이 서로를 추적하는 것을 봅니다. 이는 임금 상승이 소비자 물가 인플레이션으로 이어진다는 것을 시사합니다.

\subsubsection*{임금 인플레이션과 실업률: 필립스 곡선}

이러한 추론을 한 단계 더 나아가 봅시다. 임금은 언제 가장 많이 오를까요? 무엇이 임금 인플레이션을 유발할까요? 임금은 보통 직원들이 가장 많은 \ bargaining power를 가질 때 오릅니다. 이는 실업률이 낮고 일자리 경쟁이 적을 때입니다. 실업률이 높을 때는 고용주가 구직자를 유치하기 위해 더 높은 임금을 지불할 필요가 없습니다.

이 관계는 아서 필립스(Arthur Phillips)에 의해 처음 언급되었습니다. 그는 영국에서 실업률과 임금 상승률을 그래프로 그렸습니다. 그는 두 결과물 사이에 \textbf{음의 상관관계(negative correlation)}를 발견했습니다.

\begin{conceptbox}
\textbf{필립스 곡선 (Phillips Curve)}
원래는 실업률과 임금 상승률 간의 음의 상관관계를 나타냈지만, 오늘날에는 주로 \textbf{실업률과 소비자 물가 인플레이션 간의 상충 관계(trade-off)}를 설명하는 데 사용됩니다.
\begin{itemize}
    \item \textbf{실업률이 낮을 때}: 노동 시장이 타이트해져 기업들이 인력을 확보하기 위해 더 높은 임금을 지불해야 합니다. 이는 생산 비용을 증가시키고, 결국 소비자 물가 상승(인플레이션)으로 이어집니다.
    \item \textbf{실업률이 높을 때}: 노동 시장에 여유가 있어 기업들이 임금을 크게 올릴 필요가 없습니다. 이는 인플레이션 압력을 낮춥니다.
\end{itemize}
\end{conceptbox}

실업률이 낮을 때는 임금이 오르는 경향이 있습니다. 실업률이 높을 때는 임금이 줄어드는 경향이 있습니다. 이 관계를 종종 \textbf{필립스 곡선(Phillips Curve)}이라고 부릅니다. 오늘날에는 필립스 곡선이라는 용어가 실업률과 소비자 물가 인플레이션 간의 관계를 설명하는 데 사용됩니다. 우리는 앞서 CPI와 임금 상승률이 서로를 추적한다는 것을 보았으므로 이는 이해할 수 있습니다. 이제 최근 몇 년간 미국의 필립스 곡선을 살펴보겠습니다.

여기 1960년대 데이터가 있습니다. 이전과 마찬가지로 X축에는 실업률이 있지만, 이제 Y축에는 소비자 물가 인플레이션이 있습니다. 인플레이션과 실업률 간의 음의 관계가 유지되는 것을 볼 수 있습니다.

\subsubsection*{정책 입안자에게 미치는 영향}

왜 우리는 신경 써야 할까요? 정책 입안자에게 미치는 영향은 무엇일까요? 이 그래프는 실업률과 인플레이션 간의 명확한 상충 관계를 시사합니다. 이는 실업률을 줄이는 모든 정부 정책이 인플레이션을 증가시킬 수 있음을 의미합니다. 이것은 통화정책에 중요합니다. 물가 안정과 최대 고용이 연방준비제도의 책무라는 것을 기억하십시오. 이 목표들은 서로 상충됩니다. 두 목표를 동시에 달성하는 것은 균형을 잡는 행위입니다. 실제로 Fed의 주요 의사 결정자들이 필립스 곡선에 대해 생각하는 것을 볼 수 있습니다.

예를 들어, 2008년 당시 연방준비제도 이사회 부의장이었던 도널드 콘(Donald Cohn)은 "인플레이션의 경제 모델은 통화정책 심의에 필수적인 투입물이다. 필립스 곡선 전통의 모델은 대부분의 학술 연구자들과 정책 입안자들, 나를 포함하여, 인플레이션 변동에 대해 생각하는 방식의 핵심으로 남아 있다. 실제로 대안적인 프레임워크는 견고한 경제적 기반과 경험적 지지를 결여하고 있는 것으로 보인다"고 말했습니다.

\subsubsection*{필립스 곡선의 이동: 인플레이션 기대의 역할}

우리는 끝났을까요? 글쎄요, 아직 아닙니다. 1960년대와 1970년대 초에는 실업률과 인플레이션 간의 관계가 안정적이라고 생각되었습니다. 그러나 다음 두 십 년 동안 다른 패턴이 나타났습니다. 여기 1970년대와 1980년대 데이터가 있습니다. 불행히도 음의 상관관계는 사라진 것처럼 보입니다.

그러나 조금 자세히 보면, 일반적인 패턴이 나타나는 것 같습니다. 하나의 관계가 아니라 세 가지 관계가 있는 것처럼 보입니다. 1970년대 초, 1970년대 후반, 그리고 1980년대 초에 각각 하나씩 말입니다. 세 개의 필립스 곡선이 있는 것처럼 보입니다. 이상하죠? 이 곡선들은 시간이 지남에 따라 바깥쪽으로 이동하는 것처럼 보입니다. 바깥쪽으로 이동하는 것은 좋은 징조가 아닙니다. 필립스 곡선이 바깥쪽으로 더 많이 이동할수록 인플레이션을 억제하기 위해 더 많은 실업률이 필요합니다. 예를 들어, 인플레이션을 6\%로 유지하고 싶다고 가정해 봅시다. 1973년에는 5\%의 실업률이 필요했을 것입니다. 그러나 1982년에는 동일한 인플레이션율을 달성하기 위해 실업률이 10\%에 도달했을 것입니다.

1970년대와 1980년대는 인플레이션과 실업률 간의 전반적인 관계에 의문을 제기했지만, 데이터는 적어도 단기적으로는 필립스 곡선이 작동한다는 것을 보여줍니다. 이제 1970년대와 1980년대의 이러한 단기 필립스 곡선을 결정하는 요인을 살펴보겠습니다.

1960년에서 1995년 사이 미국의 소비자 물가 인플레이션을 살펴보겠습니다. 분명히 1970년대 초, 1970년대 후반, 1980년대 초에는 소비자 물가 인플레이션이 1960년대에 비해 더 높고 증가했습니다. 따라서 필립스 곡선의 바깥쪽 이동에 대한 그럴듯한 설명은 1970년대 초부터 1980년대 초까지 \textbf{인플레이션 기대(inflation expectations)}의 상향 이동일 수 있습니다.

\begin{conceptbox}
\textbf{인플레이션 기대 (Inflation Expectations)}
미래의 인플레이션에 대한 경제 주체들(가계, 기업, 투자자)의 예상입니다. 이 기대는 실제 인플레이션에 큰 영향을 미칩니다.
\begin{itemize}
    \item \textbf{왜 중요한가요?} 만약 사람들이 미래에 물가가 많이 오를 것이라고 예상한다면, 그들은 현재 더 높은 임금을 요구하거나, 상품 가격을 미리 올리려 할 것입니다. 이는 실제 인플레이션을 더욱 부추기는 결과를 낳습니다.
    \item \textbf{필립스 곡선에 미치는 영향}: 인플레이션 기대가 상승하면, 동일한 실업률 수준에서도 더 높은 인플레이션이 발생하게 되어 필립스 곡선이 바깥쪽(오른쪽)으로 이동합니다. 이는 정책 입안자들이 인플레이션을 억제하기 위해 더 높은 실업률을 감수해야 함을 의미합니다.
\end{itemize}
\end{conceptbox}

미래 인플레이션에 대한 기대가 왜 중요할까요? 이해하기 어렵지 않습니다. 고용주가 5\% 임금 인상을 발표했다고 가정해 봅시다. 이는 큰 인상처럼 들립니다. 그러나 소비자 물가 인플레이션이 현재 10\%로 진행 중이고 내년에도 10\%로 유지될 것으로 예상된다면 어떨까요? 이는 인플레이션을 조정한 소득, 즉 \textbf{실질 임금(real wage)}이 5\% 감소한다는 것을 의미합니다. 좋은 소식이 아닙니다. 실질 임금 인상만이 생활 수준을 높입니다. 이는 근로자들이 더 높은 인플레이션을 예상할 때 더 높은 임금을 요구하게 된다는 것을 시사합니다. 이는 임금과 실업률 간의 관계를 변화시키고, 결과적으로 실업률과 소비자 물가 인플레이션 간의 관계에 연쇄적인 영향을 미칩니다.

인플레이션 기대의 증가는 1970년대와 1980년대에 우리가 보았던 단기 필립스 곡선의 바깥쪽 이동을 설명합니다. 또한 정책 입안자들이 인플레이션 기대를 면밀히 주시하는 이유를 설명합니다. 그들은 인플레이션 기대를 \textbf{잘 고정(well anchored)}시키려고 노력합니다.

\begin{summarybox}
\textbf{레슨 3-1.2 요약}
\begin{itemize}
    \item \textbf{임금과 인플레이션의 연동}: 임금 상승은 생산 비용 증가를 통해 소비자 물가 인플레이션으로 이어지는 경향이 있습니다.
    \item \textbf{필립스 곡선}: 인플레이션과 실업률 사이에는 단기적인 상충 관계가 존재합니다. 중앙은행은 이 두 목표 사이에서 균형을 찾아야 합니다.
    \item \textbf{인플레이션 기대의 중요성}: 인플레이션 기대의 변화는 필립스 곡선을 이동시킬 수 있습니다. 기대 인플레이션이 높아지면, 동일한 실업률에서도 더 높은 인플레이션이 발생하여 필립스 곡선이 바깥쪽으로 이동합니다. 중앙은행은 인플레이션 기대를 안정적으로 유지하는 것을 중요하게 생각합니다.
\end{itemize}
\end{summarybox}

\newpage
\subsection{레슨 3-1.3: 통화정책 규칙}

안녕하세요, 통화정책 규칙 설계에 대한 수업에 오신 것을 환영합니다. 오늘날 중앙은행은 목표 달성을 위해 금리를 설정합니다. 연방준비제도의 경우, \textbf{연방 기금 금리(federal funds rate)}를 조정합니다. 그러나 Fed는 물가 안정과 최대 고용이라는 목표 사이에서 상충 관계에 직면합니다. 필립스 곡선을 기억하시나요? 이 세션에서는 이 상충 관계가 실제로 어떻게 접근되는지 논의할 것입니다. 우리는 \textbf{생산량 갭(output gap)}과 \textbf{인플레이션 갭(inflation gap)}이라는 아이디어를 소개할 것입니다. 그 다음, 중앙은행이 이러한 갭에 대해 무엇을 해야 하는지 질문할 것입니다. 우리는 통화정책이 경제 상황에 반응해서는 안 된다고 말하는 밀턴 프리드먼(Milton Friedman)의 \textbf{K-규칙(K-rule)}을 살펴볼 것입니다. 그리고 나서 금리를 인플레이션 및 고용의 갭 또는 변동과 연결하는 유명한 \textbf{테일러 준칙(Taylor Rule)}을 검토할 것입니다.

\subsubsection*{최대 고용과 물가 안정 목표 측정}

연방준비제도의 이중 책무, 즉 최대 고용과 물가 안정을 기억하십시오. 먼저 최대 고용에 초점을 맞추고 이 목표가 어떻게 달성될 수 있는지 설명해 봅시다. 실업률과 GDP는 밀접하게 연결되어 있습니다. 이 관계를 \textbf{오쿤의 법칙(Okun's Law)}이라고 합니다. 결과적으로 우리는 종종 경제 성장과 실업률을 서로 바꿔가며 이야기합니다.

경제가 자원을 완전히 활용하고 있을 때, 우리는 경제가 \textbf{잠재 GDP(potential GDP)}를 생산하고 있다고 말합니다. 경제가 잠재 GDP에 있을 때, 이는 \textbf{완전 고용(full employment)} 또는 \textbf{최대 고용(maximum employment)}을 의미합니다. 이것이 우리가 바랄 수 있는 최선입니다. 최대 고용을 달성하기 위해 중앙은행은 경제가 잠재 GDP에 가까운지 알아야 합니다.

\begin{conceptbox}
\textbf{잠재 GDP (Potential GDP)}
한 경제가 모든 생산 요소를 (노동, 자본, 토지 등) 정상적으로 활용하여 달성할 수 있는 최대 생산량 수준을 의미합니다. 이는 지속 가능한 성장을 나타내며, 이 수준에서는 인플레이션 압력이 크게 발생하지 않고 실업률은 자연 실업률 수준에 머무릅니다.

\textbf{생산량 갭 (Output Gap 또는 GDP Gap)}
실제 GDP와 잠재 GDP 간의 차이를 백분율로 나타낸 것입니다.
\begin{itemize}
    \item \textbf{양의 생산량 갭 (Positive Output Gap)}: 실제 GDP가 잠재 GDP보다 높을 때 발생합니다. 경제가 과열되어 자원을 과도하게 사용하고 있음을 의미하며, 인플레이션 압력이 높아질 수 있습니다. (예: 근로자들이 초과 근무를 하고 공장이 최대 생산 능력을 초과하여 가동될 때)
    \item \textbf{음의 생산량 갭 (Negative Output Gap)}: 실제 GDP가 잠재 GDP보다 낮을 때 발생합니다. 경제가 침체되어 자원이 충분히 활용되지 못하고 있음을 의미하며, 실업률이 높아질 수 있습니다.
\end{itemize}
중앙은행은 이 생산량 갭을 통화정책 결정의 중요한 지표로 활용하여 경제를 잠재 GDP 수준으로 되돌리려 합니다.
\end{conceptbox}

데이터에서 이것은 어떻게 보일까요? 차트에서 2030년까지 미국의 잠재 실질 GDP 추정치를 볼 수 있습니다. 미국 정부 기관인 의회 예산국(Congressional Budget Office)이 잠재 GDP를 추정합니다. 이것은 순전히 이론적인 구성물입니다. 기술, 교육 및 인구 통계학적 추세의 예상 개선으로 인해 시간이 지남에 따라 성장합니다.

여기 실제 실질 GDP가 있습니다. 두 가지가 밀접하게 추적되지만, 눈에 띄는 편차가 있습니다.

이 차트는 이러한 편차를 보여줍니다. 이는 실제 실질 GDP와 잠재 GDP 간의 퍼센트 포인트 차이로 계산됩니다. 이를 \textbf{생산량 갭(output gap)} 또는 \textbf{GDP 갭(GDP gap)}이라고 합니다. 생산량 갭이 양수인 시기가 있습니다. 이들은 경제가 과열되는 호황기입니다. 자원이 장기적으로 사용될 것보다 더 많이 사용됩니다. 근로자들은 초과 근무를 하고 추가 교대 근무를 하는 등입니다. 양의 갭은 덜 흔하고 규모도 작습니다.

대조적으로, 음의 생산량 갭은 크고 더 오래 지속될 수 있습니다.

이제 실업률을 겹쳐서 봅시다. 생산량 갭이 음수일 때 실업률이 평소보다 높다는 것을 알 수 있습니다. 2008년에 시작된 십 년은 명확한 예시입니다. 여기서 핵심 통찰력은 무엇일까요? 최대 고용이 목표 중 하나일 때 생산량 갭을 통화정책의 투입물로 사용하는 것은 매우 합리적입니다.

\subsubsection*{물가 안정 목표 측정: 인플레이션 갭}

다음으로, 연방준비제도 책무의 두 번째 부분은 물가 안정입니다. 물가 안정은 낮은 인플레이션율을 의미합니다. 하지만 얼마나 낮아야 할까요? 존 테일러(John Taylor)와 같은 학자들은 중앙은행이 2\% 목표를 사용해야 한다고 제안했습니다. 연방준비제도도 동의하는 것 같습니다. 2012년부터 연방공개시장위원회(FOMC)는 연간 2\%의 소비자 물가 인플레이션율이 연방준비제도의 법정 책무와 장기적으로 가장 일치한다고 발표했습니다. 마찬가지로 유럽중앙은행(ECB)은 2\% 미만이지만 2\%에 가까운 인플레이션을 목표로 합니다.

생산량 갭과 마찬가지로 \textbf{인플레이션 갭(inflation gap)}을 계산할 수 있습니다. 이는 인플레이션이 중앙은행의 2\% 목표율과 얼마나 다른지를 측정합니다.

\begin{conceptbox}
\textbf{인플레이션 갭 (Inflation Gap)}
실제 인플레이션율과 중앙은행이 목표로 하는 인플레이션율(예: 2\%) 간의 차이입니다.
\begin{itemize}
    \item \textbf{양의 인플레이션 갭}: 실제 인플레이션이 목표보다 높을 때 발생합니다. 중앙은행은 인플레이션을 낮추기 위해 긴축 통화정책을 고려할 수 있습니다.
    \item \textbf{음의 인플레이션 갭}: 실제 인플레이션이 목표보다 낮을 때 발생합니다. 중앙은행은 경제를 부양하고 인플레이션을 목표 수준으로 끌어올리기 위해 완화 통화정책을 고려할 수 있습니다.
\end{itemize}
\end{conceptbox}

이 그래프는 개인 소비 지출(Personal Consumption Expenditure, PCE)의 연간 성장률을 보여줍니다. 이는 소비자 물가 인플레이션의 측정치입니다. 1970년대와 1980년대에는 인플레이션이 2\% 목표보다 높았다는 것을 알 수 있습니다. 반면에 1990년대 중반부터 2020년까지는 인플레이션이 지속적으로 2\% 미만이었습니다.

\subsubsection*{통화정책 규칙: K-규칙과 테일러 준칙}

중앙은행은 실업률과 인플레이션의 변동에 어떻게 대응해야 할까요? 전혀 대응하지 않아야 할까요? 노벨상 수상 경제학자 밀턴 프리드먼(Milton Friedman)은 중앙은행이 큰 차이를 만들 수 있다는 것에 회의적이었습니다.

대신 프리드먼은 오늘날 \textbf{K-규칙(K-rule)}으로 알려진 것을 제안했습니다. 중앙은행은 단순히 매년 일정한 K 퍼센트, 예를 들어 3\%만큼 통화 공급을 증가시켜야 합니다. 분명히 이 규칙은 경제 상황을 무시합니다. 중앙은행은 재량권이 없고 미리 프로그래밍된 알고리즘을 따를 뿐입니다. 왜 이것이 좋은 규칙일까요? 통화 공급의 대규모 증가가 인플레이션을 유발한다면, 적당한 통화 성장은 물가를 안정적으로 유지해야 합니다. 이 규칙의 또 다른 장점은 금융 시장과 기업이 모든 통화정책 조치를 완벽하게 예측하고 이에 따라 결정을 내릴 수 있다는 것입니다. 통화 공급을 통제하고 예측 가능하게 유지하는 것이 경제 및 금융 안정을 촉진할 것이라는 믿음입니다.

다른 종류의 통화정책 규칙은 경제 상황에 반응하도록 설계되었습니다. 경기 순환이 비용이 많이 든다면, 이를 완화하는 것이 사회에 이로울 수 있습니다. 가장 유명한 규칙은 소위 \textbf{테일러 준칙(Taylor Rule)}으로, 금리 결정을 생산량과 인플레이션의 변동과 연결합니다.

\begin{conceptbox}
\textbf{밀턴 프리드먼의 K-규칙 (Friedman's K-rule)}
\begin{itemize}
    \item \textbf{핵심}: 중앙은행은 매년 통화 공급을 일정한 비율(K\%)로만 증가시켜야 합니다. 경제 상황에 관계없이 기계적으로 적용됩니다.
    \item \textbf{장점}: 통화 공급의 예측 가능성을 높여 인플레이션을 안정시키고, 시장의 불확실성을 줄입니다. 중앙은행의 재량적 판단으로 인한 오류를 방지합니다.
    \item \textbf{단점}: 경제 상황 변화(예: 경기 침체, 공급 충격)에 유연하게 대응할 수 없어, 경기 안정화에 실패할 수 있습니다.
\end{itemize}
\end{conceptbox}

\subsubsection*{테일러 준칙 (Taylor Rule)}

테일러 준칙을 살펴보겠습니다. 이는 1990년대에 학자 존 테일러(John Taylor)에 의해 처음 제안되었습니다. 먼저, \textbf{명목 금리(nominal interest rate)}는 \textbf{실질 금리(real interest rate, R)}와 \textbf{인플레이션($\pi$)}의 합과 같다는 것을 기억하십시오. 인플레이션이 2\% 목표에 있다고 가정해 봅시다. 또한 균형 실질 금리가 0.5\%라고 가정해 봅시다. 인플레이션이 목표에 머물러 있다고 가정하면, 명목 금리는 2.5\%가 될 것입니다. 이것이 중앙은행의 개입 없이 시장이 제공할 금리입니다. 그러나 중앙은행은 원한다면 다른 명목 금리를 선택할 수 있습니다.

사실, 중앙은행이 인플레이션 갭과 생산량 갭을 가중치 알파($\alpha$)와 베타($\beta$)로 명시적으로 고려하여 목표 금리를 선택한다고 가정해 봅시다. 이것이 여러분의 간단한 테일러 준칙입니다.

\begin{equation*}
\text{목표 금리} = \text{실질 금리} + \text{목표 인플레이션} + \alpha \times (\text{인플레이션 갭}) + \beta \times (\text{생산량 갭})
\end{equation*}

이 공식은 목표 금리가 시장이 원하는 금리에 경제 상황에 대한 조정을 더한 것이어야 한다고 말합니다.

더 정확히 말하면, 인플레이션 갭이 양수이거나 생산량 갭이 양수일 때 중앙은행은 금리를 인상해야 한다고 말합니다. 이는 경제를 식힐 것입니다.

그리고 인플레이션 갭이 음수이거나 생산량 갭이 음수일 때는 그 반대로 해야 한다고 말합니다. 이는 경제를 부양할 것입니다. 이 간단한 규칙은 연방준비제도의 의사 결정 요소를 포착합니다. 물가 안정과 최대 고용 모두 공식에 나타납니다. 물론 인플레이션 가중치 알파($\alpha$)와 생산량 갭 가중치 베타($\beta$)에 대한 값을 찾아야 합니다. 테일러의 원래 권고는 동일한 가중치를 부여하여 둘 다 0.5로 설정하는 것이었습니다. 경제 연구자들은 이후 이러한 가중치가 무엇이어야 하는지에 대한 수많은 연구를 수행했습니다. 오늘날 인플레이션 가중치 알파($\alpha$)에 대한 합의는 1.5인 반면, 생산량 갭 가중치 베타($\beta$)에 대한 추정치는 0.5에서 1 사이로 다양합니다.

\begin{examplebox}
\textbf{테일러 준칙 예시}
\begin{itemize}
    \item \textbf{기본 설정}: 실질 금리 = 0.5\%, 목표 인플레이션 = 2\%
    \item \textbf{가중치}: $\alpha = 1.5$, $\beta = 0.5$ (테일러의 원래 권고)
\end{itemize}
\textbf{시나리오 1: 경제 과열 (양의 갭)}
\begin{itemize}
    \item 실제 인플레이션 = 3\% (인플레이션 갭 = +1\%)
    \item 실제 GDP = 잠재 GDP + 2\% (생산량 갭 = +2\%)
    \item \textbf{목표 금리} = 0.5\% + 2\% + (1.5 $\times$ 1\%) + (0.5 $\times$ 2\%) \\
    = 0.5\% + 2\% + 1.5\% + 1\% = \textbf{5\%}
\end{itemize}
$\rightarrow$ 경제가 과열되어 인플레이션과 생산량이 목표보다 높으므로, 중앙은행은 금리를 5\%로 인상하여 경제를 식히려고 합니다.

\textbf{시나리오 2: 경제 침체 (음의 갭)}
\begin{itemize}
    \item 실제 인플레이션 = 1\% (인플레이션 갭 = -1\%)
    \item 실제 GDP = 잠재 GDP - 2\% (생산량 갭 = -2\%)
    \item \textbf{목표 금리} = 0.5\% + 2\% + (1.5 $\times$ -1\%) + (0.5 $\times$ -2\%) \\
    = 0.5\% + 2\% - 1.5\% - 1\% = \textbf{0\%}
\end{itemize}
$\rightarrow$ 경제가 침체되어 인플레이션과 생산량이 목표보다 낮으므로, 중앙은행은 금리를 0\%로 인하하여 경제를 부양하려고 합니다.
\end{examplebox}

통화정책을 사용하여 경기 순환을 안정화하는 것은 현대 시대 많은 중앙은행의 표준이었습니다. 하지만 이 접근 방식은 성공적이었을까요? 1980년대 이후 미국의 경험을 바탕으로 조심스럽게 "예"라고 결론 내릴 수 있습니다.

여기 1980년 이후 미국의 인플레이션, 실업률, 금리가 있습니다. 1990년 경기 침체가 시작되었을 때, 연방준비제도는 금리를 인하했습니다. 실업률은 4년 이내에 경기 침체 이전 수준 이하로 떨어졌고, 인플레이션은 1990년대 내내 낮게 유지되었습니다. 이 패턴은 2000년대 초 닷컴 버블 붕괴 이후에도 반복됩니다. 더 일반적으로, 1990년대 이후 거시경제 변수의 낮은 변동성은 중앙은행과 그들의 금리 결정에 기인합니다.

\begin{summarybox}
\textbf{레슨 3-1.3 요약}
\begin{itemize}
    \item \textbf{통화정책 목표 측정}: 최대 고용은 생산량 갭(실제 GDP와 잠재 GDP의 차이)으로, 물가 안정은 인플레이션 갭(실제 인플레이션과 목표 인플레이션의 차이)으로 측정됩니다.
    \item \textbf{통화정책 규칙}:
    \begin{itemize}
        \item \textbf{프리드먼의 K-규칙}: 경제 상황과 무관하게 통화 공급을 일정한 비율로 증가시키는 규칙입니다. 예측 가능성을 높이지만 유연성이 부족합니다.
        \item \textbf{테일러 준칙}: 금리 결정을 인플레이션 갭과 생산량 갭에 연동시키는 규칙입니다. 경제 상황에 따라 금리를 조정하여 경기 안정화를 목표로 합니다.
    \end{itemize}
    \item \textbf{성공적인 개입}: 1980년대와 2000년대 초 미국의 경험은 중앙은행의 금리 개입이 거시경제 변동성을 낮추는 데 성공적이었음을 시사합니다.
\end{itemize}
\end{summarybox}

\newpage
\section{레슨 3-2: 통화정책의 전달}

\subsection{레슨 3-2.1: 통화정책 변화가 경제에 어떻게 영향을 미치는가?}

안녕하세요, 통화정책의 경제 전달에 대한 강의에 오신 것을 환영합니다. 이 세션에서는 통화정책 결정이 경제 전반에 어떻게 확산되는지 요약할 것입니다. 미국에서 연방준비제도는 변화하는 경제 상황에 대응하여 통화정책을 조정합니다. 목표는 최대 고용과 물가 안정이라는 것을 기억하십시오. 가장 일반적인 대응은 단기 금리인 \textbf{연방 기금 금리(federal funds rate)}를 인상하거나 인하하는 것입니다.

연방 기금 금리의 변화는 다른 금리와 더 넓은 금융 시장 상황의 변화로 이어질 것입니다. 차례로, 이러한 변화는 개별 가계와 기업의 지출 결정에 영향을 미칠 것입니다. 마지막으로, 이러한 개별 효과는 거시경제적 규모로 집계될 것입니다. 이는 경제 전반의 성장, 고용 및 인플레이션으로 이어져야 합니다. 우리의 목표는 이러한 일련의 사건이 어떻게 전개되는지 이해하는 것입니다. 이는 통화정책이 어떻게 작동하는지에 대한 우리의 이해를 향상시킬 것입니다. 시작해 봅시다.

\subsubsection*{통화정책이 시장 금리에 미치는 영향}

2013년 당시 연방준비제도 이사였던 제레미 스타인(Jeremy Stein)이 연설에서 말했듯이, "통화정책은 한 가지 중요한 장점을 가지고 있는데, 그것은 모든 틈새에 스며든다는 것입니다. 상업 은행, 브로커-딜러, 역외 헤지 펀드, 특수 목적 자산 담보부 기업 어음(ABCP) 특수 목적 법인(SPV)이 공통적으로 직면하는 한 가지는 모두 동일한 시장 금리 세트를 마주한다는 것입니다. 시장 금리가 위험 선호도나 만기 변환에 참여하려는 인센티브에 영향을 미치는 한, 금리 변화는 시장의 모든 구석에 도달할 수 있습니다."

먼저 통화정책이 시장 금리에 미치는 영향을 살펴보겠습니다. 연방 기금 금리는 \textbf{초단기 금리(overnight interest rate)}입니다. 은행들은 원하면 이 금리를 벌 수 있습니다. 다른 단기 금리들은 대체 투자입니다. 이는 다른 단기 금리들이 연방 기금 금리의 변화를 추적해야 한다는 것을 의미합니다.

\begin{examplebox}
\textbf{연방 기금 금리와 단기 금리}
연방 기금 금리는 은행들이 서로에게 하루짜리 돈을 빌려줄 때 적용되는 금리입니다. 만약 Fed가 이 금리를 올리면, 은행들은 더 높은 비용으로 돈을 빌리게 되므로, 다른 단기 대출(예: 3개월 만기 국채)에 대한 금리도 따라서 오르게 됩니다. 이는 은행들이 더 높은 수익률을 요구하거나, 더 높은 비용을 대출자에게 전가하기 때문입니다.
\end{examplebox}

3개월 만기 국채(Treasury Bills) 수익률, 즉 수익률과 연방 기금 금리를 비교해 봅시다. 그래프는 둘 사이의 차이가 0에 가깝다는 것을 보여줍니다. 이는 3개월 만기 국채 수익률이 연방 기금 금리와 거의 1대1로 움직인다는 것을 의미합니다.

장기 금리는 어떨까요? 그래프는 이제 10년 만기 국채(Treasury note) 수익률과 연방 기금 금리 간의 차이를 보여줍니다. 1대1 관계가 없다는 것을 알 수 있습니다.

더 일반적으로, 연방 기금 금리 변화가 단기 금리에 미치는 영향은 클 것입니다. 그러나 장기 금리에 미치는 영향은 완화될 것입니다. 장기 금리는 더 긴 기간에 대한 기대를 포함합니다. 이러한 금리는 사람들이 미래에 연방 기금 금리가 어떻게 변할 것으로 예상하는지에 반응하며, 단지 오늘날의 변화에만 반응하는 것이 아닙니다. 예를 들어, 시장이 통화정책이 향후 몇 년 동안 상당히 긴축될 것이라고 믿는다고 가정해 봅시다. 이 경우 오늘날의 장기 금리는 단기 금리보다 훨씬 높을 수 있습니다. 이것이 미래 연방 기금 금리 변화에 대한 \textbf{선제적 발언(forward-looking statements)}이 금융 시장 참가자들에 의해 면밀히 조사되는 한 가지 이유입니다.

\subsubsection*{가계 및 기업에 미치는 영향}

다음으로, 금리 변화가 가계와 기업에 어떻게 영향을 미칠까요? 통화정책 완화로 인해 장기 금리가 하락하고 있다고 가정해 봅시다. 이는 기업이 대출을 받아 투자하는 비용을 더 저렴하게 만들 것입니다. 마찬가지로, 소비자 대출과 모기지(mortgages)도 더 저렴해집니다. 이는 가전제품과 같은 내구재 및 주택 소유에 대한 지출을 자극합니다.

장기 금리 하락은 또한 미국 달러 자산의 수익률을 국제 투자자들에게 덜 매력적으로 만들 것입니다.

이러한 투자자들은 달러 표시 자산에 대한 수요를 줄여 달러 가치를 낮출 수 있습니다. \textbf{약한 달러(weaker dollar)}는 미국 상품이 해외에서 더 저렴해지므로 미국 수출을 증가시키는 경향이 있습니다. 동시에 수입품은 더 비싸져 미국 소비자와 기업이 국내 생산품을 구매하도록 유도합니다.

이 모든 조각들을 합쳐 봅시다. 기업과 가계에 대한 영향이 있습니다. 장기 금리 변화는 주가에도 영향을 미칩니다. 일반적으로 우리는 기업의 시장 가치를 모든 미래 배당금의 할인된 가치로 생각합니다. 장기 금리가 하락하면 미래 배당금에 대한 할인율이 더 유리해집니다. 이는 기업의 시장 가치를 증가시키고 결과적으로 주가를 높입니다. 이러한 주가 상승은 개인의 부를 증가시키고, 이는 다시 소비를 더욱 증가시킵니다.

\begin{examplebox}
\textbf{금리 인하의 경제적 파급 효과}
Fed가 금리를 인하하면:
\begin{enumerate}
    \item \textbf{대출 비용 감소}: 기업은 더 저렴하게 자금을 조달하여 투자를 늘리고, 가계는 모기지나 자동차 대출을 더 낮은 이자로 받을 수 있어 소비를 늘립니다.
    \item \textbf{주가 상승}: 미래 기업 이익의 현재 가치가 높아져 주가가 상승하고, 이는 가계의 자산 증가로 이어져 소비를 더욱 촉진합니다.
    \item \textbf{환율 하락}: 달러 자산의 매력이 감소하여 달러 가치가 하락합니다. 이는 수출을 늘리고 수입을 줄여 국내 생산을 활성화합니다.
\end{enumerate}
이러한 효과들이 종합적으로 작용하여 경제 성장, 고용 증가, 그리고 필립스 곡선에 따라 인플레이션 상승으로 이어질 수 있습니다.
\end{examplebox}

낮은 금리로 경제를 부양하는 것이 실업률을 낮출 것이라는 믿음입니다. 이는 또한 필립스 곡선에 따라 물가를 높일 것입니다. 이것이 이론입니다. 이제 이러한 효과가 데이터에서 어떻게 나타나는지 살펴보겠습니다. 우리는 연방 기금 금리의 예상치 못한 인상을 검토할 것입니다. 산업 생산, 통화 공급, 실업률, 인플레이션에 어떤 일이 일어나는지 알고 싶습니다. 우리가 보게 될 차트들은 월별 데이터를 기반으로 추정된 시계열 통계 분석에 기반합니다. 우리는 2년의 시간 범위에 걸쳐 효과를 측정할 것입니다.

\subsubsection*{데이터로 본 통화정책의 영향}

첫 번째 그래프는 연방 기금 금리가 50bp(베이시스 포인트) 예상치 못하게 한 번 인상된 후 어떻게 되는지 보여줍니다. 단기적으로는 금리가 더 인상될 가능성이 있다는 점에 유의하십시오. 그러나 시간이 지남에 따라 금리 인상의 효과는 사라지고 우리는 시작점으로 돌아갑니다.

여기 M2 통화 공급이 있습니다. 연방준비제도는 공개 시장 조작(open market operations)을 통해 연방 기금 금리 인상을 시행합니다. 공개 시장 조작은 Fed가 증권을 매각하고 그 대가로 은행 준비금 형태로 현금을 받는다는 것을 기억하십시오. 준비금은 M2로 측정되는 통화 공급의 일부입니다. 연방 기금 금리를 인상하면 통화 공급이 줄어듭니다. 이 효과가 가라앉는 데 약 2년이 걸립니다.

여기 산업 생산이 있습니다. 그래프는 산업 생산이 감소하는 데 약 6개월이 걸린다는 것을 보여줍니다. 이것은 중요한 점입니다. 통화정책은 시차를 두고 경제에 영향을 미칩니다. 이는 적어도 두 가지 이유로 발생합니다. 첫째, 기업은 이미 투자 계획을 위한 자금을 받았을 수 있으며 여전히 이를 실행하고 있을 수 있습니다. 둘째, 금리 변화가 소비자 지출에 영향을 미치는 데 시간이 걸릴 수 있습니다.

실업률은 산업 생산의 거울 이미지입니다. 몇 달이 걸리지만, 산업 생산이 감소함에 따라 실업률은 상승하기 시작합니다. 또한 연방 기금 금리가 인상되면 인플레이션이 하락한다는 증거가 있습니다. 중앙은행이 경제 과열을 우려하여 연방 기금 금리를 인상하면, 중기적으로 실업률은 증가하고 인플레이션은 하락할 것입니다.

\begin{summarybox}
\textbf{레슨 3-2.1 요약}
\begin{itemize}
    \item \textbf{통화정책 전달 경로}: 연방 기금 금리 변화 $\rightarrow$ 다른 금리 및 금융 시장 상황 변화 $\rightarrow$ 가계 및 기업의 지출 결정 변화 $\rightarrow$ 거시경제적 성장, 고용, 인플레이션 변화.
    \item \textbf{시장 금리 영향}: 연방 기금 금리는 단기 금리에 큰 영향을 미치지만, 장기 금리에는 미래 기대가 반영되어 영향이 완화됩니다.
    \item \textbf{경제적 영향}: 금리 인하(완화 정책)는 기업 투자 및 가계 소비를 자극하고, 달러 약세를 유발하여 수출을 증대시키며, 주가를 상승시켜 부를 증가시킵니다.
    \item \textbf{데이터로 본 효과}: 연방 기금 금리 인상 시, 통화 공급은 즉시 감소하고, 산업 생산은 약 6개월 후 감소하며, 실업률은 상승하고, 인플레이션은 하락합니다. 이러한 효과는 \textbf{시차(time lag)}를 두고 나타납니다.
\end{itemize}
\end{summarybox}

\newpage
\subsection{레슨 3-2.2: 은행을 통한 통화정책 전달}

안녕하세요, 통화정책의 은행 대출 채널에 대한 강의에 오신 것을 환영합니다. 이 수업에서는 통화정책이 은행의 신용 공급에 어떻게 영향을 미치는지 살펴볼 것입니다. 우리는 세 가지 메커니즘을 논의할 것입니다: \textbf{순이자 마진(net interest rate margin)}, \textbf{지급 준비금(reserves)}, 그리고 \textbf{예금(deposits)}.

\subsubsection*{은행의 사업 모델과 순이자 마진}

은행의 사업 모델을 생각해 봅시다. 은행은 예금자와 머니 마켓 펀드로부터 단기적으로 자금을 빌려, 이 자금을 사용하여 장기 대출을 실행합니다. 단기 금리가 보통 장기 금리보다 낮기 때문에, 은행은 단기 금리와 장기 금리 간의 차이를 이용하여 돈을 법니다.

\begin{conceptbox}
\textbf{순이자 마진 (Net Interest Margin, NIM)}
은행의 핵심 수익성 지표입니다. 은행이 대출 및 모기지에서 얻는 이자 수입과 예금자에게 지급하는 이자 비용 간의 차이를 나타냅니다.
\begin{itemize}
    \item \textbf{계산}: (대출 이자 수입 - 예금 이자 비용) / 총 자산
    \item \textbf{의미}: NIM이 높을수록 은행의 수익성이 좋다는 것을 의미하며, 이는 새로운 대출을 실행할 수 있는 자본이 더 많다는 것을 뜻합니다.
\end{itemize}
\end{conceptbox}

은행의 대출 및 모기지에서 발생하는 순이자 수입과 은행이 예금자에게 지급하는 이자 비용 간의 차이를 \textbf{순이자 마진(net interest margin)}이라고 합니다. 순이자 마진은 은행의 수익성을 측정합니다. 순이자 마진이 클수록 이익이 높습니다.

이 그래프는 1980년대 이후 미국 은행의 순이자 마진을 보여줍니다. 보시다시피, 은행의 순이자 마진은 시간이 지남에 따라 감소했습니다.

동시에, 2010년경 연방준비제도가 연방 기금 금리를 처음으로 0으로 인하했을 때 순이자 마진이 크게 증가하는 것을 볼 수 있습니다.

\begin{examplebox}
\textbf{통화정책이 순이자 마진에 미치는 영향}
\begin{itemize}
    \item \textbf{시나리오 1: 금리 인상 (긴축 정책)}
    \begin{itemize}
        \item 은행이 100달러를 2\% 이자로 예금받아, 5\% 이자로 100달러 고정 금리 모기지를 대출했다고 가정합니다. 순이자 마진은 3달러(3\%)입니다.
        \item Fed가 연방 기금 금리를 2\%에서 3\%로 인상하면, 은행은 예금 이자를 3\%로 올려야 예금을 유지할 수 있습니다.
        \item 하지만 고정 금리 모기지 이자는 5\%로 변하지 않습니다.
        \item \textbf{결과}: 순이자 마진은 5달러 - 3달러 = 2달러(2\%)로 감소합니다.
    \end{itemize}
    \item \textbf{시나리오 2: 금리 인하 (완화 정책)}
    \begin{itemize}
        \item 위와 동일한 초기 조건에서 Fed가 금리를 인하하면, 은행은 예금 이자를 낮출 수 있습니다.
        \item 고정 금리 대출의 이자는 변하지 않으므로, 순이자 마진은 증가하여 은행의 수익성이 개선됩니다.
    \end{itemize}
\end{itemize}
\textbf{결론}: 금리 인상은 은행의 순이자 마진을 압박하여 수익을 감소시키고, 이는 새로운 대출을 위한 자본을 줄여 신용 공급을 위축시킵니다. 반대로 금리 인하는 순이자 마진을 개선하여 신용 공급을 확대할 수 있습니다.
\end{examplebox}

통화정책이 순이자 마진에 어떻게 영향을 미치는지 간단한 예시로 살펴보겠습니다. 100달러의 예금으로 전액 자금을 조달하는 은행을 생각해 봅시다. 예금 이자율은 2\%이므로 총 이자 비용은 2달러입니다.

은행은 5\% 이자를 지급하는 100달러 고정 금리 모기지 하나를 발행했습니다.

이자 수입은 5달러이고, 2달러의 이자 비용을 제외하면 순이자 마진은 3달러 또는 3\%입니다.

이제 연방준비제도가 연방 기금 금리를 2\%에서 3\%로 인상한다고 가정해 봅시다. 연방 기금 시장은 초단기 대출 시장이므로, 다른 모든 초단기 및 단기 금리도 인상될 것입니다. 예금은 요구불 예금이며, 따라서 초단기 대출과 매우 유사합니다. 은행이 예금을 잃고 싶지 않다면, 그에 따라 예금 금리를 인상해야 할 것입니다. 은행은 예금에 대해 3\%를 지불해야 하지만, 고정 금리 모기지의 이자율은 변하지 않았습니다. 따라서 순이자 마진은 3\%에서 2\%로 하락합니다.

순이자 마진 하락은 이익 감소를 의미합니다. 이익 감소는 다시 새로운 대출을 위한 자본 감소를 의미합니다. 다시 말해, 통화정책은 은행의 이익에 직접적으로 영향을 미치고, 따라서 은행의 신용 공급에 영향을 미칩니다. 물론 통화정책이 이익에 얼마나 강하게 영향을 미치는지는 은행이 대차대조표에 보유하고 있는 고정 금리 대출의 비중에 따라 달라집니다. 기업 대출과 같은 일부 대출은 일반적으로 변동 금리 대출이므로, 이러한 대출의 순이자 마진은 통화정책의 영향을 받지 않습니다.

\subsubsection*{지급 준비금 채널}

통화정책이 은행에 영향을 미치는 두 번째 방법은 은행이 연방준비제도 계좌에 보유하고 있는 \textbf{지급 준비금(reserves)}을 통해서입니다. 이를 \textbf{연방 기금(federal funds)}이라고 합니다. 연방준비제도가 연방 기금 금리를 인상할 때, 일반적으로 공개 시장 조작을 통해 그렇게 합니다. 즉, 연방준비제도는 은행에 증권을 매각합니다. 은행은 준비금으로 증권 구매 대금을 지불합니다. 연방 기금 시장에서 사용할 수 있는 준비금이 적어지면 단기 자금 시장의 신용 공급이 감소합니다. 은행은 단기 유동성 필요, 예를 들어 대출을 발행할 때 연방 기금 시장을 사용합니다. 물론 은행은 다른 시장에서도 자금을 빌릴 수 있지만, 종종 그 자금은 더 비싸거나 환매 조건부 채권(Repo) 시장과 같이 담보를 요구합니다. 은행에 사용할 수 있는 자금이 적거나, 남아 있는 자금이 더 비쌀 때, 은행은 신용 공급을 줄입니다. 따라서 통화정책은 준비금을 줄임으로써 은행의 신용 공급을 줄입니다.

\subsubsection*{예금 채널}

통화정책이 은행에 영향을 미치는 또 다른 방법은 \textbf{예금(deposits)}을 통해서입니다. 순이자 마진에 대해 제가 들었던 간단한 예시를 기억하십시오. 그 예시에서 은행은 연방 기금 금리 인상에 맞춰 예금 금리를 인상했습니다. 그러나 여러분 대부분이 당좌 예금 및 저축 계좌에서 알고 있듯이, 은행은 일반적으로 당좌 예금 및 저축 계좌의 금리 인상에 빠르게 대응하지 않습니다.

대부분의 예금자들은 투자 대안을 찾지 않는 경향이 있습니다. 시간이 걸리기 때문입니다. 모든 은행 정보를 변경해야 하므로 은행을 바꾸는 데는 훨씬 더 많은 시간이 걸립니다. 그럼에도 불구하고 일부 예금자, 특히 상당한 예금을 가진 예금자들은 은행을 바꿉니다. 그들은 예금 금리와 대체 투자 금리 간의 차이가 증가할 때 은행에서 예금을 인출하는 경향이 있습니다. 예금의 더 가까운 대체재는 \textbf{머니 마켓 뮤추얼 펀드(money market mutual funds)}입니다.

그래프는 머니 마켓 뮤추얼 펀드의 총 자산과 연방 기금 금리를 보여줍니다. 보시다시피, 연방준비제도가 연방 기금 금리를 인상할 때, 예를 들어 2004년에서 2006년 사이 또는 2016년에서 2019년 사이에 머니 마켓 뮤추얼 펀드의 자산이 훨씬 더 빠르게 증가했습니다.

예금 감소는 또한 은행이 자금을 덜 보유하고 있으므로 신용 공급을 줄여야 한다는 것을 의미합니다. 여러분은 왜 이러한 행동이 은행에 최적인지 궁금할 수 있습니다. 은행은 일부 예금을 잃고 일부 대출 기회를 포기하는 것과 모든 예금의 이자율을 1\% 포인트 인상하고 현재 대출 속도를 유지하는 것 사이에서 선택에 직면합니다. 다르게 말하면, 예금의 5\%를 잃고 대출을 5\% 줄이는 것과 모든 예금의 이자율을 1\% 포인트 인상하고 현재 대출 속도를 유지하는 것 중 어느 것을 선호하시겠습니까? 최근 연구에 따르면 은행은 예금 금리를 인상하는 시점에 대해 전략적입니다. 미국 내에서는 은행 경쟁이 더 심한 곳에서 예금 금리가 더 빠르게 인상되는 경향이 있습니다. 이는 은행을 바꾸기 쉬울수록 은행이 예금을 유지하기 위해 금리 변화에 더 많이 반응해야 한다는 것을 시사합니다. 이 발견의 반대 측면은 은행 시스템이 고도로 집중된 지역이나 국가에서는 이러한 통화정책의 예금 채널이 덜 중요하다는 것입니다.

\begin{summarybox}
\textbf{레슨 3-2.2 요약: 은행을 통한 통화정책 전달 3가지 채널}
\begin{itemize}
    \item \textbf{순이자 마진 (Net Interest Margin) 채널}: 연방 기금 금리 인상은 은행의 순이자 마진을 압박하여 수익을 감소시키고, 이는 새로운 대출을 위한 자본을 줄여 신용 공급을 위축시킵니다.
    \item \textbf{지급 준비금 (Reserves) 채널}: 연방 기금 금리 인상은 공개 시장 조작을 통해 은행 시스템 내의 지급 준비금을 감소시킵니다. 이는 은행의 단기 자금 조달 비용을 높여 신용 공급을 줄입니다.
    \item \textbf{예금 (Deposits) 채널}: 연방 기금 금리 인상은 예금자들이 더 높은 수익률을 찾아 머니 마켓 뮤추얼 펀드와 같은 대체 투자로 이동하게 만듭니다. 이는 은행의 예금 기반 자금을 감소시켜 신용 공급을 줄입니다.
\end{itemize}
이 세 가지 채널 모두 연방 기금 금리 인상이 은행의 전반적인 신용 공급을 감소시킨다는 것을 의미합니다.
\end{summarybox}

\newpage
\subsection{레슨 3-2.3: 기업과 가계를 통한 통화정책 전달}

안녕하세요, 기업과 가계를 통한 통화정책 전달에 대한 강의에 오신 것을 환영합니다. 이 수업의 목표는 금리 변화가 차입 비용으로 어떻게 전달되는지 더 잘 이해하는 것입니다. 우리는 이것이 기업, 특히 변동 금리 은행 대출에 의존하는 소규모 및 재정적으로 제약된 기업에 어떻게 중요한지 살펴볼 것입니다. 이는 경제의 투자 및 일자리에 연쇄적인 영향을 미 미칠 수 있습니다. 그 다음, 동일한 채널이 가계에 어떻게 영향을 미치는지 검토할 것입니다. 여기서는 주택 시장과 모기지 대출에 초점을 맞출 것입니다. 우리는 금리 변화가 모기지를 재융자할 옵션이 있는 가계에 어떻게 중요한지 살펴볼 것입니다. 그리고 나서 경제의 총수요에 대한 함의를 논의할 것입니다. 시작해 봅시다.

\subsubsection*{기업에 미치는 영향: 변동 금리 대출}

먼저 기업을 살펴보겠습니다. 거의 모든 기업은 어떤 형태의 부채를 안고 있습니다. 가장 큰 기업들은 고정 금리 부채인 채권 시장에 접근할 수 있습니다. 소규모 기업들은 일반적으로 은행 대출에만 접근할 수 있습니다. 이러한 사업 대출은 종종 \textbf{변동 금리(floating or variable interest rate)}를 가집니다. 이는 대출 이자율이 일부 변동 벤치마크에 고정 마크업을 더한 것과 같다는 것을 의미합니다. 변동 벤치마크는 과거에는 런던 은행 간 금리(LIBOR), 담보부 초단기 자금 조달 금리(SOFR) 또는 최우수 고객에게 제공되는 금리인 프라임 금리(prime rate) 중 하나일 수 있습니다. 일반적으로 이는 기업이 지불하는 이자율이 단기 금리와 함께 움직인다는 것을 의미합니다.

\begin{conceptbox}
\textbf{변동 금리 대출 (Variable Rate Loans)}
대출 이자율이 시장 금리(벤치마크 금리)에 따라 변동하는 대출입니다.
\begin{itemize}
    \item \textbf{구성}: 벤치마크 금리 (예: SOFR, Prime Rate) + 고정 마크업 (은행의 이윤)
    \item \textbf{특징}: 중앙은행의 정책 금리 변화에 직접적으로 영향을 받습니다. 정책 금리가 오르면 대출 이자도 오르고, 내리면 대출 이자도 내립니다.
    \item \textbf{주요 대상}: 주로 소규모 기업 대출이나 일부 가계 대출(예: 변동 금리 모기지)에 적용됩니다. 대기업은 고정 금리 채권 발행을 통해 금리 변동 위험을 회피하는 경향이 있습니다.
\end{itemize}
\end{conceptbox}

이것이 통화정책과 어떻게 연결되는지 예시를 통해 살펴보겠습니다. 한 기업이 100만 달러의 사업 대출을 가지고 있다고 가정해 봅시다. 이 대출의 이자율은 SOFR + 2\%라고 가정해 봅시다. 연방 기금 금리와 마찬가지로 SOFR은 초단기 금리이며 둘은 밀접하게 움직입니다. 연방 기금 금리와 SOFR이 현재 4\%라고 가정해 봅시다. 이는 대출 이자율이 6\%이고 총 연간 이자 지급액이 60,000달러라는 것을 의미합니다.

이제 FOMC가 연방 기금 금리를 1\%로 인하하기로 결정했다고 가정해 봅시다. SOFR도 이 하락 움직임을 따를 것입니다. 이는 기업 대출의 이자율이 6\%에서 3\%로 재설정되어 이자 지급액이 30,000달러로 감소한다는 것을 의미합니다. 따라서 연방 기금 금리 인하는 기업이 투자 및 고용에 사용할 수 있는 30,000달러를 확보해 줍니다.

더 일반적으로, 연방 기금 금리 인하는 레버리지가 있는 기업의 이익을 증가시킬 것입니다. 더 높은 이익의 또 다른 함의는 기업의 대차대조표와 주식 가치가 개선될 것이라는 점입니다. 더 높은 이익으로 기업은 더 나은 재정 상태에 있습니다. 이는 차입 비용을 더욱 줄일 수 있습니다.

데이터가 이러한 아이디어를 지지할까요? 글쎄요, 그렇기도 하고 아니기도 합니다. 연구는 기업이 연방 기금 금리 인상에 어떻게 반응하는지, 그리고 그것이 경제 전체에 중요한지 여부를 조사했습니다. 연구에 따르면 대기업은 금리가 인상될 때 대부분 영향을 받지 않습니다. 이자 비용이 증가하고 수익이 감소하더라도 생산 및 고용 수준을 유지합니다. 이러한 대기업은 기업 어음 시장 및 기타 단기 신용원에 접근하는 데 문제가 없습니다. 이는 이익이 감소하더라도 단기 차입을 늘릴 수 있도록 합니다. 동시에 대기업의 재고는 증가합니다.

이와는 대조적으로, 통화정책이 소규모 기업의 차입 비용에 미치는 영향은 훨씬 더 중요합니다. 소규모의 재정적으로 제약된 기업은 이자 비용 증가에 재고를 줄이고 근무 시간과 생산을 줄임으로써 반응합니다. 경제 전체에 중요한 것은, 그 효과가 여전히 상당하다는 것입니다. 이는 소규모 기업이 여전히 미국 전체 근로자의 절반 이상을 고용하고 있기 때문입니다.

\subsubsection*{가계에 미치는 영향: 주택 시장과 재융자}

다음으로, 통화정책에 가계가 어떻게 반응하는지 살펴보겠습니다. 모든 가계가 FOMC의 성명을 주의 깊게 읽는다고 말하는 것은 아닙니다. 그러나 그들의 의사 결정은 변화하는 시장 금리를 포함한 지배적인 경제 상황에 반응할 것입니다. 가계 대차대조표에서 가장 큰 항목은 보통 \textbf{모기지(mortgage)}입니다.

이 그래프는 30년 고정 금리 모기지 이자율과 연방 기금 금리를 보여줍니다. 보시다시피, 두 금리는 함께 움직이므로 모기지 금리는 통화정책 변화에 반응할 것입니다.

모기지 금리 변화에서 얼마나 큰 효과를 기대할 수 있을까요? 답은 가계가 얼마나 부채를 지고 있는지에 달려 있습니다. 가계가 부채가 전혀 없다면, 차입 비용 변화는 무관할 것입니다. 지도는 2007년, 2008년 금융 위기 직전 미국 모든 카운티의 가계 중위 부채 대비 소득 비율을 보여줍니다. 보시다시피, 많은 카운티에는 연간 소득의 세 배가 넘는 부채를 가진 가계가 있습니다. 연간 소득 대비 이렇게 높은 차입으로 인해 금리 변화는 가계 재정에 상당한 영향을 미칠 수 있습니다.

미국의 일부 모기지는 변동 금리이지만, 대부분의 모기지는 고정 금리입니다. FOMC가 연방 기금 금리를 인하하기로 결정하면 모기지 금리도 하락하는 경향이 있습니다. 이에 대응하여 가계는 현재 모기지를 \textbf{재융자(refinance)}하기로 결정할 수 있습니다. 이는 가계가 현재 모기지를 더 낮은 이자율 모기지로 대체하는 것을 의미합니다. 또한 이는 선불 비용을 수반할 수 있습니다. 재융자로 인한 이자율 절감은 상당할 수 있습니다. 모기지 이자율 그래프로 돌아가서 21세기에 초점을 맞춰 봅시다.

2006년 모기지 이자율은 약 6.5\%였고, 2011년 12월에는 4\% 미만으로 떨어졌습니다. 원금을 전혀 갚지 않은 가계는 모기지를 재융자함으로써 모기지 이자 비용을 거의 40\% 절감할 수 있었고, 많은 사람들이 그렇게 했습니다. 최근 연구에 따르면 재융자를 한 평균 가계는 월 모기지 상환액에서 거의 1,000달러를 절약했습니다.

가계는 이 추가 현금으로 무엇을 했을까요? 모기지 상환액을 절약한 가계는 더 많이 소비하는 경향이 있습니다. 예를 들어, 사람들은 외식을 하거나 집을 개조하기 시작할 수 있습니다. 내구재와 자동차에 대한 지출이 증가합니다. 결과적으로 상품 및 서비스에 대한 수요가 증가하여 기업에 이익이 되고 경제를 활성화합니다.

\begin{summarybox}
\textbf{레슨 3-2.3 요약: 기업과 가계를 통한 통화정책 전달}
\begin{itemize}
    \item \textbf{기업에 미치는 영향}:
    \begin{itemize}
        \item \textbf{변동 금리 대출}: 소규모 기업은 주로 변동 금리 은행 대출에 의존하므로, 정책 금리 인하 시 이자 비용이 감소하여 이익이 증가하고 투자 및 고용을 늘릴 수 있습니다.
        \item \textbf{대기업과의 차이}: 대기업은 채권 시장 접근성 등으로 인해 금리 변화에 덜 민감합니다.
    \end{itemize}
    \item \textbf{가계에 미치는 영향}:
    \begin{itemize}
        \item \textbf{모기지 금리}: 정책 금리 변화는 모기지 금리에 영향을 미칩니다.
        \item \textbf{재융자}: 금리 인하 시 가계는 모기지를 재융자하여 월 상환액을 줄일 수 있으며, 이로 인해 가처분 소득이 증가하여 소비(내구재, 자동차 등)를 늘립니다.
    \end{itemize}
    \item \textbf{총수요 증가}: 금리 인하로 인한 기업의 투자 증가와 가계의 소비 증가는 경제의 전반적인 상품 및 서비스 수요를 증가시켜 경제를 활성화합니다.
\end{itemize}
\end{summarybox}

\newpage
\section{레슨 3-3: 통화정책의 한계}

\subsection{레슨 3-3.1: 통화정책의 한계: 공급 충격}

안녕하세요, 스태그플레이션에 대한 강의에 오신 것을 환영합니다. 이 수업에서는 낮은 성장, 높은 인플레이션, 높은 실업률로 특징지어졌던 1970년대 기간을 살펴볼 것입니다. 이 사례 연구는 통화정책의 상충 관계와 한계를 더 잘 이해하는 데 도움이 될 것입니다. \textbf{스태그플레이션(stagflation)}이라는 용어는 \textbf{경기 침체(stagnation)}와 \textbf{인플레이션(inflation)}의 조합입니다. 스태그플레이션은 중앙은행에게 특히 어려운 문제입니다. 경제를 부양하기 위해 중앙은행은 금리를 인하하지만, 낮은 금리는 또한 인플레이션을 증가시키는 경향이 있습니다. 성장이 낮고 인플레이션이 높은 상황에서 중앙은행은 선택의 폭이 제한적입니다.

\begin{conceptbox}
\textbf{스태그플레이션 (Stagflation)}
경제 성장(stagnation)이 둔화되거나 침체되는 동시에 물가 상승(inflation)이 지속되는 현상을 말합니다. 일반적으로 경제학 이론에서는 경기 침체 시 물가가 하락하거나 안정되는 경향이 있다고 보지만, 스태그플레이션은 이러한 일반적인 관계를 벗어나는 예외적인 상황입니다.
\begin{itemize}
    \item \textbf{중앙은행의 딜레마}:
    \begin{itemize}
        \item \textbf{경기 부양 (금리 인하)}: 실업률을 낮추고 성장을 촉진할 수 있지만, 이미 높은 인플레이션을 더욱 악화시킬 수 있습니다.
        \item \textbf{인플레이션 억제 (금리 인상)}: 인플레이션을 낮출 수 있지만, 이미 낮은 성장을 더욱 둔화시키고 실업률을 높일 수 있습니다.
    \end{itemize}
    이러한 딜레마 때문에 스태그플레이션은 중앙은행에게 매우 다루기 어려운 문제입니다.
\end{itemize}
\end{conceptbox}

그래프는 1960년부터 1990년까지의 소비자 물가 인플레이션을 보여줍니다. 1960년대 후반에 인플레이션이 가속화되기 시작한 것을 볼 수 있습니다. 가장 큰 급등은 1974년과 1980년에 발생했으며, 이때 인플레이션율은 각각 11\%와 13.5\%에 달했습니다. 이러한 급등의 원인을 살펴보겠습니다.

\subsubsection*{공급 충격: 1970년대 오일 쇼크}

주요 원인은 두 번의 큰 \textbf{오일 쇼크(oil shocks)}입니다. 첫 번째 오일 쇼크는 욤 키푸르 전쟁(Yom Kippur War) 이후 OPEC이 미국과 다른 서방 국가들에 대한 석유 금수 조치를 취하면서 발생했습니다. 몇 달 만에 유가는 배럴당 3달러에서 약 11달러로 거의 300\% 상승했습니다.

\begin{conceptbox}
\textbf{공급 충격 (Supply Shocks)}
생산 비용에 갑작스럽고 예상치 못한 변화를 가져오는 사건을 말합니다. 이는 주로 생산 요소의 가격이나 공급량에 영향을 미칩니다.
\begin{itemize}
    \item \textbf{긍정적 공급 충격}: 생산 비용을 낮추고 생산량을 늘리는 충격 (예: 기술 혁신, 풍작).
    \item \textbf{부정적 공급 충격}: 생산 비용을 높이고 생산량을 줄이는 충격 (예: 유가 급등, 자연재해, 전쟁).
\end{itemize}
\textbf{스태그플레이션과의 관계}: 부정적 공급 충격은 생산 비용을 높여 물가를 상승시키고(인플레이션), 동시에 생산량을 감소시켜 경제 성장을 둔화시키거나 침체시킬 수 있습니다(경기 침체). 이는 스태그플레이션을 유발하는 주요 원인 중 하나입니다.
\end{itemize}
\end{conceptbox}

이는 경제에 엄청난 비용 충격이었습니다. 거의 모든 경제 부문이 휘발유 형태나 발전용, 화학 산업 등에서 석유를 수입으로 사용했기 때문입니다. 생산자 물가는 급격히 상승했습니다. 동시에 소비자들은 휘발유와 난방에 더 많은 돈을 지불해야 했고, 이는 다른 소비재에 대한 지출을 줄였습니다.

불과 몇 년 후, 두 번째 오일 쇼크가 경제를 강타했습니다. 1979년 이란 혁명과 1980년 이란-이라크 전쟁 발발 이후 석유 생산량이 감소했습니다. 석유 공급 감소와 미래 석유 공급에 대한 불확실성으로 인해 유가는 폭등했습니다. 1980년에는 유가가 배럴당 39달러 50센트에 달했습니다. 이는 1973년 초 이후 10배 이상 증가한 수치입니다. 미국에서는 자동차 제조업체들이 더 연료 효율적인 자동차를 제공하기 위해 고군분투하면서 특히 큰 타격을 입었습니다. 이러한 대규모 충격은 경제 둔화로 이어졌습니다.

그래프에서 실업률을 볼 수 있습니다. 1975년에는 9\%로 제2차 세계 대전 이후 최고치를 기록했으며, 1981년 두 번째 오일 쇼크와 함께 다시 감소하여 7.8\%에 도달했습니다. 연방준비제도의 책무는 물가 안정과 최대 고용이라는 것을 기억하십시오. 그러나 오일 쇼크는 그 반대 효과, 즉 높은 인플레이션과 높은 실업률을 초래했습니다. 오일 쇼크는 \textbf{공급 충격(supply shocks)}의 예시입니다.

그래프는 공급 충격이 실업률과 인플레이션 간의 상충 관계, 즉 필립스 곡선을 어떻게 바깥쪽으로 이동시키는지 보여줍니다. 우리는 별도의 강의에서 필립스 곡선을 자세히 논의했습니다.

\subsubsection*{통화정책의 한계와 대응}

공급 충격에 직면했을 때 통화정책의 선택지는 제한적입니다. 낮은 금리는 공급 제약에 영향을 미치지 않기 때문입니다. 실제로 오일 쇼크에 대한 많은 대응은 통화정책이 아니라 단기 및 장기적으로 석유 수요에 영향을 미쳤습니다. 1973년 오일 쇼크에 대한 가장 극적인 단기 대응은 미국에서 휘발유 배급제를 시행한 것이었습니다. 예를 들어, 홀수 번호판을 가진 차량은 홀수 날에만 휘발유를 살 수 있었고, 다른 차량은 짝수 날에만 살 수 있었습니다.

석유 수요에 상당한 장기적 영향을 미친 대응은 결국 유가를 낮추는 데 도움이 되었습니다. 몇 가지 예시가 있습니다. 휘발유 소비를 줄이기 위해 1974년에는 전국 최고 속도 제한이 시속 55마일로 부과되었습니다. 1년 후인 1975년에는 기업 평균 연비(CAFE) 기준이 만들어졌습니다. 이 기준은 미국에서 자동차와 경트럭의 연비 개선을 요구했습니다. 동시에 미국은 \textbf{전략 비축유(Strategic Petroleum Reserve)}를 개발했습니다.

이러한 조치와 생산량 증가는 결국 1980년대에 소위 \textbf{오일 과잉(oil glut)}으로 이어졌습니다. 유가는 상당히 하락했고 인플레이션을 낮추는 데 기여했습니다.

1980년대의 이러한 공급 충격은 인플레이션과 실업률 간의 상충 관계를 다시 안쪽으로 이동시키는 데 도움이 되었습니다. 그래프에서 볼 수 있듯이, 1980년대 후반에는 1980년대 초에 비해 실업률을 줄이기 위해 훨씬 적은 인플레이션이 필요했습니다.

그럼에도 불구하고 통화정책은 공급 충격에 대응하는 데 거의 할 수 없습니다. 이 기간은 인플레이션 기대의 변화를 야기했습니다. 즉, 합의는 낮은 인플레이션 체제에서 인플레이션이 높게 유지될 체제로 바뀌었습니다.

높은 인플레이션 기대는 높은 임금 요구로 이어졌고, 이는 다시 인플레이션을 증가시켰습니다.

폴 볼커(Paul Volcker)가 1979년 연방준비제도 의장이 되었을 때, 그는 인플레이션과 인플레이션 기대를 낮추기를 원했습니다.

그는 연방 기금 금리를 인상함으로써 그렇게 했으며, 1981년 12월에는 19\%에 달했습니다. 금리의 급격한 인상은 유가 상승과 함께 1980년과 1982년 경기 침체를 촉발했습니다.

실업률은 10\% 이상으로 상승하여 제2차 세계 대전 종전부터 코로나19 팬데믹까지 최고 수준을 기록했습니다. 다시 말해, 이 기간 동안 연방준비제도는 물가 안정과 최대 고용이라는 이중 책무에도 불구하고 실업률보다는 인플레이션에 초점을 맞추기로 선택했습니다. 연방 기금 금리 인상은 인플레이션의 큰 하락으로 이어졌습니다.

더 중요하게는, 연방준비제도는 인플레이션을 낮추기 위해 일시적으로 더 높은 실업률이라는 대가를 받아들였습니다. 이는 미래의 낮은 인플레이션에 대한 신뢰할 수 있는 약속으로 여겨졌습니다. 이 약속은 다시 인플레이션 기대를 낮추는 데 도움이 되었고, 이는 낮은 인플레이션 체제를 확립하는 데 기여했습니다.

\begin{summarybox}
\textbf{레슨 3-3.1 요약: 통화정책의 한계와 공급 충격}
\begin{itemize}
    \item \textbf{스태그플레이션}: 1970년대는 낮은 성장, 높은 인플레이션, 높은 실업률이 동시에 나타나는 스태그플레이션으로 특징지어졌습니다. 이는 중앙은행에게 큰 딜레마를 안겨줍니다.
    \item \textbf{공급 충격의 영향}: 유가 급등과 같은 공급 충격은 생산 비용을 증가시켜 물가를 상승시키고(인플레이션), 동시에 경제 활동을 위축시켜 실업률을 높입니다. 이는 필립스 곡선을 바깥쪽으로 이동시킵니다.
    \item \textbf{통화정책의 한계}: 통화정책은 금리 조정을 통해 총수요에 영향을 미치지만, 공급 측면의 제약(예: 원유 생산량 감소)을 직접적으로 해결할 수 없습니다. 따라서 공급 충격에 대응하는 데는 한계가 있습니다.
    \item \textbf{인플레이션 기대 관리}: 공급 충격이 인플레이션 기대를 영구적으로 변화시킬 경우, 중앙은행은 인플레이션과 인플레이션 기대를 낮추기 위해 일시적으로 높은 실업률을 감수하는 강력한 긴축 정책을 펼칠 수 있습니다 (예: 폴 볼커 의장 시기). 이는 미래의 물가 안정에 대한 신뢰를 구축하는 데 중요합니다.
\end{itemize}
\end{summarybox}

\newpage
\subsection{레슨 3-3.2: 통화정책의 한계: 고용 없는 회복과 위험 감수 인센티브}

(이 레슨의 본문 내용은 제공된 텍스트 파트 1/2에 포함되어 있지 않습니다.)

\newpage
\section*{체크리스트: 학습 및 복습 가이드}

이 문서를 통해 학습한 내용을 점검하고 복습하는 데 도움이 되는 체크리스트입니다.

\begin{itemize}
    \item \textbf{개요 및 기본 개념 이해}
    \begin{itemize}
        \item Fed의 이중 책무(최대 고용, 물가 안정)를 정확히 설명할 수 있는가?
        \item GDP, 실질 GDP, 실업률의 정의와 중요성을 이해하는가?
    \end{itemize}
    \item \textbf{인플레이션과 실업률의 관계}
    \begin{itemize}
        \item 오쿤의 법칙(실업률과 GDP 성장률의 음의 상관관계)을 설명할 수 있는가?
        \item 자연 실업률의 개념과 중앙은행이 실업률을 0으로 만들지 않는 이유를 아는가?
        \item 소비자 물가 인플레이션과 생산자 물가 인플레이션, 임금 인플레이션의 관계를 이해하는가?
        \item 필립스 곡선(인플레이션과 실업률의 상충 관계)을 설명할 수 있는가?
        \item 인플레이션 기대가 필립스 곡선에 미치는 영향(곡선 이동)을 이해하는가?
    \end{itemize}
    \item \textbf{통화정책 규칙}
    \begin{itemize}
        \item 잠재 GDP와 생산량 갭(output gap)의 개념을 아는가?
        \item 인플레이션 갭의 개념을 아는가?
        \item 밀턴 프리드먼의 K-규칙과 테일러 준칙의 차이점을 설명할 수 있는가?
        \item 테일러 준칙의 구성 요소(실질 금리, 목표 인플레이션, 인플레이션 갭, 생산량 갭)를 이해하는가?
        \item 테일러 준칙에 따른 금리 결정 예시를 설명할 수 있는가?
    \end{itemize}
    \item \textbf{통화정책 전달 메커니즘}
    \begin{itemize}
        \item 연방 기금 금리 변화가 단기 및 장기 시장 금리에 미치는 영향을 설명할 수 있는가?
        \item 금리 변화가 환율, 주가, 소비, 투자에 미치는 영향을 설명할 수 있는가?
        \item 통화정책이 경제에 시차를 두고 영향을 미치는 이유를 아는가?
        \item 은행을 통한 통화정책 전달의 세 가지 채널(순이자 마진, 지급 준비금, 예금)을 설명할 수 있는가?
        \item 기업(특히 소규모 기업의 변동 금리 대출)과 가계(모기지 재융자)를 통한 통화정책 전달을 설명할 수 있는가?
    \end{itemize}
    \item \textbf{통화정책의 한계}
    \begin{itemize}
        \item 스태그플레이션의 개념과 중앙은행의 딜레마를 설명할 수 있는가?
        \item 공급 충격(예: 오일 쇼크)이 스태그플레이션을 유발하는 과정을 설명할 수 있는가?
        \item 통화정책이 공급 충격에 효과적으로 대응하기 어려운 이유를 아는가?
        \item 인플레이션 기대를 안정화하기 위한 중앙은행의 역할과 그 중요성을 이해하는가?
    \end{itemize}
\end{itemize}

\newpage
\section*{FAQ: 자주 묻는 질문}

초심자들이 혼동할 수 있는 질문들을 Q\&A 형식으로 정리했습니다.

\begin{itemize}
    \item \textbf{Q: Fed의 이중 책무인 '최대 고용'과 '물가 안정'은 항상 함께 달성될 수 있나요?}
    \textbf{A:} 이상적으로는 그렇지만, 현실에서는 종종 상충 관계에 놓입니다. 예를 들어, 경제를 과열시켜 고용을 늘리려 하면 인플레이션이 높아질 수 있고, 반대로 인플레이션을 잡으려 하면 고용이 위축될 수 있습니다. Fed는 이 두 목표 사이에서 균형을 찾아야 합니다.

    \item \textbf{Q: 오쿤의 법칙에서 실업률 1\% 감소가 GDP 2\% 증가로 이어진다는 것은 항상 정확한가요?}
    \textbf{A:} 오쿤의 법칙은 경험적 관계이며, '약 -2'라는 기울기는 역사적 데이터에서 관찰된 평균적인 경향입니다. 특정 시기나 국가에 따라 이 관계는 달라질 수 있으며, 정확히 2배가 아닌 근사치로 이해하는 것이 중요합니다.

    \item \textbf{Q: 필립스 곡선이 '사라졌다'는 말은 무슨 의미인가요?}
    \textbf{A:} 1970년대와 1980년대에 높은 인플레이션과 높은 실업률이 동시에 나타나는 스태그플레이션이 발생하면서, 기존의 안정적인 필립스 곡선 관계가 깨진 것처럼 보였습니다. 이는 인플레이션 기대의 변화로 인해 필립스 곡선 자체가 이동했기 때문이며, 단기적인 상충 관계는 여전히 유효하다고 해석됩니다.

    \item \textbf{Q: 테일러 준칙은 중앙은행이 반드시 따라야 하는 법적 규칙인가요?}
    \textbf{A:} 아닙니다. 테일러 준칙은 학자들이 제안한 통화정책의 '지침' 또는 '규칙' 모델입니다. 중앙은행이 금리 결정을 할 때 참고할 수 있는 유용한 프레임워크를 제공하지만, 법적 구속력은 없으며, 실제 정책 결정은 다양한 요소를 종합적으로 고려하여 이루어집니다.

    \item \textbf{Q: 통화정책이 경제에 영향을 미치는 데 왜 '시차'가 발생하나요?}
    \textbf{A:} 금리 변화가 기업의 투자 결정이나 가계의 소비 결정에 영향을 미치기까지는 시간이 걸립니다. 기업은 이미 투자 계획을 세웠을 수 있고, 가계는 대출 금리 변화를 인지하고 행동을 바꾸기까지 시간이 필요하기 때문입니다. 이러한 시차 때문에 중앙은행은 미래 경제 상황을 예측하여 선제적으로 정책을 결정해야 합니다.

    \item \textbf{Q: 공급 충격이 발생했을 때 통화정책은 왜 효과적이지 않나요?}
    \textbf{A:} 통화정책은 주로 총수요를 조절하여 경제에 영향을 미칩니다. 하지만 공급 충격(예: 유가 급등)은 생산 비용을 직접적으로 증가시켜 물가를 올리고 생산량을 줄입니다. 금리 인하로 수요를 늘려도 생산 비용 문제는 해결되지 않으며, 오히려 인플레이션을 악화시킬 수 있습니다. 따라서 공급 충격에는 통화정책보다 공급 측면의 구조적 정책이 더 효과적일 수 있습니다.
\end{itemize}

\newpage
\section*{빠르게 훑어보기: 1페이지 요약}

\begin{tcolorbox}[colback=myblue!10!white, colframe=myblue!50!black, title=\textbf{모듈 3 핵심 요약: 통화정책의 경제적 영향}]
\textbf{1. Fed의 이중 책무:}
\begin{itemize}
    \item \textbf{최대 고용}: 자연 실업률 수준의 고용 달성.
    \item \textbf{물가 안정}: 낮은 인플레이션 유지 (목표 2\%).
\end{itemize}

\textbf{2. 인플레이션과 실업률의 관계:}
\begin{itemize}
    \item \textbf{오쿤의 법칙}: 실업률 $\uparrow$ $\leftrightarrow$ GDP 성장률 $\downarrow$ (음의 상관관계).
    \item \textbf{필립스 곡선}: 인플레이션 $\leftrightarrow$ 실업률 (단기적 상충 관계).
    \item \textbf{인플레이션 기대}: 필립스 곡선을 이동시키는 핵심 요인. 기대 $\uparrow$ $\rightarrow$ 곡선 바깥 이동.
\end{itemize}

\textbf{3. 통화정책 규칙:}
\begin{itemize}
    \item \textbf{생산량 갭}: 실제 GDP와 잠재 GDP의 차이.
    \item \textbf{인플레이션 갭}: 실제 인플레이션과 목표 인플레이션의 차이.
    \item \textbf{테일러 준칙}: 금리 = 실질금리 + 목표인플 + $\alpha$(인플 갭) + $\beta$(생산량 갭). 경제 상황에 따라 금리 조정.
\end{itemize}

\textbf{4. 통화정책 전달 메커니즘:}
\begin{itemize}
    \item \textbf{금리 채널}: 정책 금리 $\rightarrow$ 시장 금리 $\rightarrow$ 투자/소비/환율/주가.
    \item \textbf{은행 대출 채널}: 정책 금리 $\rightarrow$ 순이자 마진, 지급 준비금, 예금 $\rightarrow$ 은행 신용 공급.
    \item \textbf{기업/가계 채널}: 정책 금리 $\rightarrow$ 변동 금리 대출 비용, 모기지 재융자 $\rightarrow$ 기업 투자, 가계 소비.
    \item \textbf{시차}: 정책 효과는 경제에 시차를 두고 나타남.
\end{itemize}

\textbf{5. 통화정책의 한계:}
\begin{itemize}
    \item \textbf{공급 충격 (예: 오일 쇼크)}: 스태그플레이션 유발. 통화정책은 수요 조절에 강하고 공급 충격 대응에 약함.
    \item \textbf{인플레이션 기대 관리}: 공급 충격 시, 인플레이션 기대를 안정화하기 위해 일시적 고실업 감수 필요 (볼커 의장 사례).
\end{itemize}
\end{tcolorbox}

\newpage

\newpage
\section*{개요: 통화 정책의 한계}

\begin{summarybox}{문서 전체 핵심 요약}
이 문서는 통화 정책이 효과적으로 작동하지 않는 두 가지 주요 상황을 탐구합니다.
첫째, \textbf{고용 없는 회복(Jobless Recovery)} 현상으로, 경기 침체 후 GDP는 성장하지만 노동 시장은 느리게 회복되는 경우입니다.
이는 높은 불확실성이나 경제의 구조적 변화로 인해 발생할 수 있으며, 통화 정책만으로는 해결하기 어렵습니다.
둘째, \textbf{완화적 통화 정책의 부작용}으로, 장기간 낮은 금리가 과도한 위험 감수를 유발하여 금융 불안정이나 자산 거품을 초래할 수 있습니다.
특히 2000년대 초반의 주택 시장 거품 사례를 통해 이러한 위험을 상세히 분석합니다.
\end{summarybox}

\newpage
\section{통화 정책의 한계: 고용 없는 회복과 위험 감수 유인}

\subsection{개요}
이번 강의에서는 통화 정책이 기대만큼 효과를 발휘하지 못했던 시기들을 살펴봅니다. 특히 두 가지 주요 관련 문제에 집중할 것입니다. 첫째는 '고용 없는 회복(Jobless Recovery)'이라는 개념입니다. 이 용어는 1991년 경기 침체 이후의 회복을 설명하기 위해 처음 사용되었습니다. 연방준비제도(Federal Reserve)의 상당한 금리 인하에도 불구하고, 노동 시장은 예상보다 느리게 회복되었습니다. 우리는 이러한 현상이 왜 발생하는지 이해하려고 노력할 것입니다.

둘째는 완화적 통화 정책의 의도치 않은 결과인 '과도한 위험 감수(Excessive Risk-Taking)'입니다. 금리 인하는 경제를 활성화할 수 있지만, 만약 금리가 너무 오랫동안 너무 낮게 유지된다면 어떻게 될까요? 일부에서는 이것이 금융 불안정을 초래할 수 있다고 주장합니다. 우리는 이 주장을 자세히 검토할 것입니다.

\subsection{고용 없는 회복 (Jobless Recoveries)}

\subsubsection{개념 및 역사적 배경}
역사적으로 경기 침체가 끝난 후에는 실업률이 급격히 하락하는 경향이 있었습니다. 예를 들어, 1982년 경기 침체 이후의 상황을 보면 이를 명확히 알 수 있습니다. 그러나 1991년 경기 침체 이후에는 대조적으로 실업률이 점진적으로만 하락했습니다. 이는 연방준비제도가 3년 동안 금리를 낮게 유지하고 1992년에서 1995년 사이에 GDP가 5\% 이상 성장했음에도 불구하고 발생한 현상입니다. 이러한 노동 시장의 느린 회복은 경제학자들을 당혹스럽게 했습니다.

\begin{examplebox}{1991년 경기 침체 후 고용 없는 회복}
1991년 경기 침체 이후, 연방준비제도는 경제 활성화를 위해 금리를 크게 인하했습니다. 일반적으로 금리 인하는 기업의 투자와 고용을 촉진하여 실업률을 빠르게 낮추는 효과가 있습니다. 실제로 1992년부터 1995년까지 미국의 GDP는 5\% 이상 성장하며 경제는 회복세를 보였습니다. 하지만 이러한 경제 성장에도 불구하고, 실업률은 과거의 경기 회복기처럼 급격히 떨어지지 않고 매우 점진적으로 하락했습니다. 이는 경제가 성장해도 일자리가 충분히 늘어나지 않는 '고용 없는 회복'의 전형적인 예시로 기록되었습니다.
\end{examplebox}

2001년 경기 침체 이후에도 동일한 패턴이 나타났습니다. 연방준비제도는 경기 침체에 대응하여 연방기금금리(federal funds rate)를 6.5\%에서 1\%로 대폭 인하했으며, 2004년 여름까지 연방기금금리를 인상하지 않았습니다. 그러나 실업률은 2003년 여름이 되어서야 하락하기 시작했고, 그마저도 하락세는 완만했습니다. 이처럼 낮은 금리 수준에서의 통화 정책 완화는 전례 없는 일이었습니다. 이러한 이유로 경제학자들은 2000년대 초반을 '고용 없는 회복' 시기로 지칭합니다.

\subsubsection{2001년 고용 없는 회복의 특징}
2001년 고용 없는 회복 시기의 노동 시장 데이터를 통해 어떤 특징이 나타났는지 살펴보겠습니다. 노동 시장 지표들은 근로자들이 일자리를 찾는 데 상당한 어려움을 겪었음을 시사했습니다.

\begin{itemize}
    \item \textbf{노동 시장 참여율 감소:} 특히 젊은 근로자들 사이에서 노동 시장 참여율이 감소했습니다. 이는 일자리를 찾으려는 의지를 가진 사람 자체가 줄어들었음을 의미할 수 있습니다.
    \item \textbf{실업 기간 증가:} 실업 상태가 지속되는 기간이 상당히 길어졌습니다. 이는 새로운 일자리를 찾는 데 더 많은 시간이 소요되었음을 보여줍니다.
    \item \textbf{구인 광고 지수 하락:} 일자리 공고의 척도인 구인 광고 지수가 경기 침체 이전 수준보다 낮게 유지되었습니다. 이는 기업들이 예상만큼 적극적으로 채용하지 않았다는 증거입니다.
\end{itemize}
이러한 증거들은 기업들이 예상만큼 많이 고용하지 않았음을 시사합니다.

또한, 2003년 9월 말 현재의 고용 통계 조사(Current Employment Statistics survey)에 따르면, 2001년 3월 경기 침체 시작 이후 비농업 부문 고용이 약 280만 개 감소했습니다. 이 중 제조업 부문이 가장 큰 타격을 입어, 전체 감소한 280만 개 일자리 중 약 240만 개가 제조업에서 발생했습니다.

\subsubsection{느린 노동 시장 회복의 원인}
그렇다면 왜 노동 시장의 반응은 그렇게 느렸을까요? 그리고 왜 제조업 부문이 가장 큰 타격을 입었을까요? 여기에는 두 가지 주요 설명이 있습니다.

\begin{enumerate}
    \item \textbf{경제에 대한 높은 불확실성 (High Uncertainty):}
    기업들은 경제 상황에 대한 불확실성이 높을 때 투자나 고용을 꺼리는 경향이 있습니다. 2000년대 초반에는 9.11 테러 공격과 그 여파(아프가니스탄 및 이라크 전쟁 포함)가 가장 명백한 불확실성의 원천이었습니다. 이러한 외부 충격은 기업들이 미래를 예측하기 어렵게 만들었고, 이는 고용 결정에 부정적인 영향을 미쳤습니다.

    \item \textbf{미국 경제의 구조적 변화 (Structural Changes):}
    제조업 고용 손실은 미국 경제의 구조적 변화 때문이었습니다. 구조적 변화는 특정 산업의 영구적인 쇠퇴를 의미합니다. 제조업의 쇠퇴는 새로운 현상이 아니었습니다. 1980년대 초반부터 제조업 부문 고용 근로자 비중은 급격히 감소하기 시작했습니다. 이는 부분적으로 국제 경쟁 심화에 의해 주도되었습니다.

    \begin{examplebox}{구조적 변화의 예시: 제조업 쇠퇴}
    1980년대 오일 위기 이후, 일본 및 서유럽 자동차 제조업체들이 미국 시장으로 확장하면서 자동차 시장의 경쟁이 심화된 것이 가장 두드러진 사례입니다. 이러한 국제 경쟁은 미국 제조업의 일자리 감소에 기여했습니다. 구조적 변화는 일반적으로 느리게 진행되는 과정이므로, 정책 입안자들이 실시간으로 그 결과를 측정하기 어렵게 만듭니다.

    2000년대 초반의 논의를 되돌아보면, 한 가지 요인이 눈에 띄게 빠져 있습니다. 2001년 12월 11일, 중국이 세계무역기구(WTO)에 가입했습니다. 그 결과, 미국의 무역 정책은 중국산 수입품에 대한 잠재적인 관세 인상을 철회했습니다. 최근 연구에 따르면, 제조업 일자리 손실은 중국과의 수입 경쟁에 직면한 미국 기업들에 집중되었습니다. 이는 중국의 WTO 가입이 미국 제조업 일자리 감소에 상당한 영향을 미쳤음을 시사합니다.
    \end{examplebox}
\end{enumerate}

\subsubsection{통화 정책의 한계점}
불확실성이 높거나 구조적 변화가 진행 중일 때 통화 정책은 효과적일 수 있을까요?

\begin{itemize}
    \item \textbf{불확실성 감소:} 일반적으로 통화 정책은 현재와 미래에 완화적인 통화 정책을 약속함으로써 미래 경제 발전에 대한 불확실성을 줄일 수 있습니다. 이는 '선제적 안내(forward guidance)'의 사용을 포함합니다. 그러나 2000년대 초반의 경우, 불확실성은 테러 공격과 그에 따른 전쟁으로 인해 발생했습니다. 이러한 유형의 불확실성은 통화 정책으로 줄이기 어려웠을 것입니다.

    \item \textbf{구조적 변화 대응:} 마찬가지로 통화 정책은 구조적 변화로 인한 고용 효과를 줄이는 데 거의 할 수 있는 것이 없습니다. 쇠퇴하는 산업에서 사라진 일자리는 결국 성장하는 산업(예: 서비스 또는 기술 산업)의 일자리로 대체될 것으로 예상됩니다. 그러나 이 과정은 시간이 걸릴 수 있습니다. 새로운 일자리가 창출되어야 하고, 근로자들은 재배치되고 재교육을 받아야 합니다.

    \item \textbf{최적의 작동 조건:} 통화 정책은 경기 순환적 변동(cyclical fluctuation)에 대응할 때 가장 잘 작동합니다. 구조적 변화가 실업률을 주도할 때는 할 수 있는 것이 거의 없습니다.
\end{itemize}

\begin{cautionbox}{통화 정책의 한계}
통화 정책은 경제의 단기적인 경기 변동(예: 경기 침체와 회복)을 완화하는 데 효과적입니다. 하지만 9.11 테러와 같은 외부 충격으로 인한 불확실성이나, 제조업 쇠퇴와 같은 장기적인 산업 구조 변화로 인한 실업 문제에는 직접적인 해결책을 제공하기 어렵습니다. 이러한 문제들은 통화 정책보다는 재정 정책, 교육 및 훈련 프로그램, 산업 정책 등 다른 정책 수단이 더 효과적일 수 있습니다.
\end{cautionbox}

\newpage
\subsection{완화적 통화 정책의 위험: 과도한 위험 감수}
이제 이 세션의 나머지 부분에서는 완화적 통화 정책의 위험을 살펴보겠습니다. 2001년, 투자 운용사 PIMCO의 경제학자 폴 맥컬리(Paul McCully)는 연방준비제도가 설정한 극도로 낮은 연방기금금리가 주택 가격 상승을 유발할 수 있다고 예측했습니다.

\subsubsection{주택 시장 거품 형성 과정}
맥컬리의 예측은 현실이 되었습니다. 다음은 주택 시장 거품이 어떻게 형성되었는지에 대한 설명입니다.

\begin{enumerate}
    \item \textbf{모기지 금리 하락:}
    2000년 8\% 이상이던 모기지 금리는 2003년 여름 5\%로 떨어졌습니다.

    \item \textbf{주택 수요 및 가격 상승:}
    모기지 금리가 낮아지면 사람들이 새 집이나 더 큰 집을 구매할 여력이 생깁니다. 이는 주택 수요를 증가시키고 주택 가격을 상승시킵니다. 실제로 2000년에서 2006년 사이에 주택 가격은 60\% 이상 상승했습니다.

    \item \textbf{건설 경기 호황 및 부의 효과:}
    주택 가격 상승의 결과로 건설 경기가 호황을 누렸습니다. 또한, '부의 효과(wealth effect)'로 인해 주택 소유자들은 모기지 재융자(mortgage refinancing)와 주택 담보 대출(home equity lines of credit, HELOC)을 이용하여 주택을 담보로 더 많은 돈을 빌렸습니다. 이렇게 얻은 자금은 자동차나 가전제품과 같은 소비재에 지출되었습니다.

    \item \textbf{신용 공급의 극적인 증가:}
    주택 담보 대출(HELOC)은 2000년 약 1,500억 달러에서 2005년 말까지 6,500억 달러로 극적으로 증가했습니다. 이러한 신용 공급의 증가는 주택 건설, 소비자 지출 및 일자리 창출을 촉진했습니다.
\end{enumerate}

\begin{examplebox}{모기지 금리 하락과 주택 시장의 변화}
\textbf{2000년:} 모기지 금리 8\% 이상. 주택 구매 부담이 컸음.
\textbf{2003년 여름:} 모기지 금리 5\%로 하락.
\begin{itemize}
    \item \textbf{구매력 증가:} 같은 월 상환액으로 더 비싼 집을 살 수 있게 됨.
    \item \textbf{수요 증가:} 더 많은 사람이 주택 시장에 진입하거나 더 큰 집으로 옮기려 함.
    \item \textbf{가격 상승:} 수요가 공급을 초과하면서 주택 가격이 2000년 대비 2006년까지 60\% 이상 폭등.
\end{itemize}
\textbf{부의 효과:} 주택 가격이 오르자, 주택 소유자들은 자신이 부자가 되었다고 느끼고, 주택 담보 대출을 통해 현금을 인출하여 소비를 늘렸습니다. 이는 경제 활성화에 기여했지만, 동시에 과도한 부채와 자산 거품의 위험을 키웠습니다.
\end{examplebox}

\subsubsection{금융 불안정 경고}
2005년, 당시 IMF 수석 경제학자 라구람 라잔(Raghuram Rajan)은 주택 거품에 대한 우려를 표명했습니다. 그는 낮은 금리와 금융 혁신이 금융 시장 참여자들이 점점 더 많은 위험을 감수하는 상황을 만들었다고 주장했습니다. 그는 이러한 위험이 금융 시스템의 안정성을 위협할 수 있다고 결론지었습니다. 돌이켜보면, 그의 분석은 정확했습니다.

\subsubsection{일반적인 원칙 및 중앙은행의 역할}
기억해야 할 더 일반적인 요점은 낮은 금리가 위험 축적(buildup of risk)으로 이어질 수 있다는 것입니다. 낮은 금리는 과도한 차입을 장려할 수 있으며, 이는 주택 시장 및 금융 시장에서 자산 가격 거품을 부채질할 수 있습니다. 낮은 금리와 위험 감수 사이의 연관성은 미국뿐만 아니라 모든 주요 경제권에서 나타났습니다.

위험 감수를 장려하는 것은 경제를 활성화하는 한 가지 방법입니다. 그러나 중앙은행은 완화적 통화 정책이 금융 시스템의 전반적인 건전성에 미치는 영향을 염두에 두어야 합니다. 2008년 금융 위기 이후, 금융 안정성(financial stability)은 연방준비제도의 임무에 추가되었습니다. 연방준비제도는 이제 경제 내의 위험 축적을 모니터링합니다.

\begin{cautionbox}{과도한 위험 감수의 위험}
낮은 금리는 기업과 가계가 돈을 빌리는 비용을 낮춰 투자를 늘리고 소비를 촉진하는 긍정적인 효과가 있습니다. 하지만 금리가 너무 오랫동안 너무 낮게 유지되면, 사람들은 더 높은 수익을 추구하기 위해 과도한 위험을 감수하게 됩니다. 이는 주식, 부동산 등 자산 시장에 거품을 형성하고, 부채 수준을 위험하게 높여 금융 시스템 전체의 안정성을 해칠 수 있습니다. 2008년 글로벌 금융 위기는 이러한 위험이 현실화된 대표적인 사례입니다.
\end{cautionbox}

\newpage
\subsection{결론 및 요약}
우리는 통화 정책의 두 가지 잠재적 한계를 논의했습니다.

\begin{enumerate}
    \item \textbf{구조적 변화 또는 외부 사건으로 인한 불확실성:}
    통화 정책은 구조적 변화나 외부 사건에 대한 불확실성이 실업률을 증가시킬 때 할 수 있는 것이 거의 없습니다. 이러한 문제들은 통화 정책의 범위를 넘어섭니다.

    \item \textbf{완화적 통화 정책의 부작용:}
    2000년대 초반의 고용 없는 회복 이후의 완화적 통화 정책은 미국에서 주택 가격 거품을 부채질했을 수 있습니다. 이 경험은 더 일반적인 문제를 강조합니다. 낮은 금리는 위험 감수를 장려하고 경제 내에 위험을 축적시킬 수 있습니다.
\end{enumerate}

\newpage
\section{용어 정리}

\begin{table}[h!]
\centering
\caption{주요 용어 정리}
\label{tab:glossary}
\begin{adjustbox}{width=\textwidth,center}
\begin{tabular}{lllc}
\toprule
\textbf{용어} & \textbf{쉬운 설명} & \textbf{원어} & \textbf{비고} \\
\midrule
고용 없는 회복 & 경제는 성장하지만 일자리는 느리게 늘어나는 현상 & Jobless Recovery & 1991년, 2001년 침체 후 관찰 \\
연방준비제도 & 미국의 중앙은행 & Federal Reserve & 통화 정책을 결정 \\
연방기금금리 & 은행 간 초단기 대출 금리 목표치 & Federal Funds Rate & 중앙은행의 주요 정책 금리 \\
완화적 통화 정책 & 금리를 낮춰 경제 활동을 촉진하는 정책 & Loose/Accommodative Monetary Policy & 경기 부양 목적 \\
과도한 위험 감수 & 낮은 금리로 인해 투자자들이 지나치게 위험한 투자를 하는 경향 & Excessive Risk-Taking & 금융 불안정의 원인 \\
구조적 변화 & 특정 산업의 영구적인 쇠퇴나 산업 구조의 변화 & Structural Change & 제조업 쇠퇴 등이 해당 \\
선제적 안내 & 중앙은행이 미래 통화 정책 방향을 미리 알려주는 것 & Forward Guidance & 시장의 불확실성 감소 목적 \\
부의 효과 & 자산 가치 상승으로 소비가 증가하는 현상 & Wealth Effect & 주택 가격 상승 시 발생 \\
모기지 재융자 & 기존 주택 담보 대출을 새로운 조건으로 다시 받는 것 & Mortgage Refinancing & 낮은 금리 활용 \\
주택 담보 대출 & 주택을 담보로 돈을 빌리는 신용 한도 & Home Equity Line of Credit (HELOC) & 주택 가치 상승 시 활용 증가 \\
금융 안정성 & 금융 시스템이 충격에 강하고 원활하게 작동하는 상태 & Financial Stability & 중앙은행의 중요한 목표 \\
\bottomrule
\end{tabular}
\end{adjustbox}
\end{table}

\newpage
\section{빠르게 훑어보기: 핵심 요약}

\begin{summarybox}{통화 정책의 한계: 1페이지 요약}
\textbf{1. 고용 없는 회복 (Jobless Recovery)}
\begin{itemize}
    \item \textbf{정의:} GDP는 성장하지만 실업률은 느리게 하락하는 현상.
    \item \textbf{사례:} 1991년 및 2001년 미국 경기 침체 후 관찰.
    \item \textbf{원인:}
    \begin{itemize}
        \item \textbf{높은 불확실성:} 9.11 테러, 전쟁 등 외부 충격.
        \item \textbf{구조적 변화:} 제조업 쇠퇴, 국제 경쟁(중국 WTO 가입).
    \end{itemize}
    \item \textbf{통화 정책의 한계:} 경기 순환적 문제에는 효과적이나, 외부 충격으로 인한 불확실성이나 구조적 실업에는 대응하기 어려움.
\end{itemize}

\textbf{2. 완화적 통화 정책의 위험: 과도한 위험 감수}
\begin{itemize}
    \item \textbf{문제점:} 금리가 너무 낮게, 너무 오랫동안 유지될 때 발생.
    \item \textbf{메커니즘:}
    \begin{itemize}
        \item 낮은 모기지 금리 $\rightarrow$ 주택 수요 증가 $\rightarrow$ 주택 가격 급등 (2000-2006년 60\% 이상 상승).
        \item 부의 효과 $\rightarrow$ 주택 담보 대출(HELOC) 증가 (2000년 1,500억 달러 $\rightarrow$ 2005년 6,500억 달러).
        \item 과도한 차입 및 자산 거품 형성.
    \end{itemize}
    \item \textbf{경고:} 라구람 라잔(2005년)은 낮은 금리와 금융 혁신이 금융 불안정을 초래할 수 있다고 경고.
    \item \textbf{중앙은행의 역할:} 2008년 금융 위기 이후, 금융 안정성 감시 및 위험 축적 모니터링이 연방준비제도의 주요 임무에 추가됨.
\end{itemize}
\end{summarybox}

\end{document}
